\documentclass[11pt,letterpaper]{article}
\usepackage[T1]{fontenc}
\usepackage{lmodern}
\usepackage[margin=1in]{geometry}
\usepackage{amsmath,amssymb,amsthm,mathtools}
\usepackage{microtype,booktabs,longtable,array,enumitem}
\usepackage{xcolor,fancyhdr,xurl,hyperref,listings}
\definecolor{ink}{RGB}{26,48,71}
\definecolor{blue}{RGB}{22,76,119}
\hypersetup{colorlinks=true,linkcolor=blue,citecolor=blue,urlcolor=blue,
 pdftitle={Distribution of zeta zeros in intuitionistic primitive recursive arithmetic},
 pdfauthor={ChatGPT, prepared for Vladimir Reshetnikov}}
\pagestyle{fancy}\fancyhf{}
\fancyhead[L]{\small Zeta zeros in iPRA}
\fancyhead[R]{\small Quantitative constructive proofs}
\fancyfoot[C]{\thepage}
\setlength{\headheight}{14pt}
\setlength{\parindent}{0pt}\setlength{\parskip}{.35em}
\setlist{itemsep=.25em,topsep=.4em}
\allowdisplaybreaks[2]\numberwithin{equation}{section}
\newtheorem{theorem}{Theorem}[section]
\newtheorem{lemma}[theorem]{Lemma}
\newtheorem{proposition}[theorem]{Proposition}
\newtheorem{corollary}[theorem]{Corollary}
\theoremstyle{definition}
\newtheorem{definition}[theorem]{Definition}
\newtheorem{remark}[theorem]{Remark}
\newcommand{\ipra}{\mathsf{iPRA}}
\newcommand{\PR}{\mathsf{PR}}
\newcommand{\N}{\mathbb N}\newcommand{\Z}{\mathbb Z}
\newcommand{\Q}{\mathbb Q}\newcommand{\R}{\mathbb R}\newcommand{\C}{\mathbb C}
\newcommand{\ii}{\mathrm i}\newcommand{\ee}{\mathrm e}
\newcommand{\normone}[1]{\lVert #1\rVert_1}
\newcommand{\abs}[1]{\lvert #1\rvert}
\newcommand{\code}[1]{\texttt{\detokenize{#1}}}
\DeclareMathOperator{\wind}{wind}
\DeclareMathOperator{\Log}{Log}
\DeclareMathOperator{\Ree}{Re}\DeclareMathOperator{\Imm}{Im}
\lstset{basicstyle=\small\ttfamily,breaklines=true,columns=fullflexible,
 keepspaces=true,frame=single,rulecolor=\color{black!20},showstringspaces=false}

\begin{document}
\hypersetup{pageanchor=false}
\begin{titlepage}
\vspace*{.55in}
{\color{ink}\Huge\bfseries Distribution of the zeros\\[.15em]of the Riemann zeta function\par}
\vspace{.35in}
{\LARGE in intuitionistic primitive recursive arithmetic\par}
\vspace{.35in}
{\large Quantitative human proofs, finite arithmetic certificates,\\
and a Lean formalization roadmap\par}
\vspace{.25in}
{\large Prepared for Vladimir Reshetnikov\par}
{\large 15 November 2026\par}
\vspace{.25in}
\begin{abstract}
We give constructive versions of the classical location, symmetry and
counting theorems for zeta zeros. The main distribution estimate is the
Riemann--von Mangoldt formula with an explicit logarithmic error. A
quantitative minimum-modulus argument supplies a primitive recursive
choice of regular contour heights. Thus the zero count is a usable
finite object, rather than a count conditional on an unspecified
zero-free boundary.

Additional results include simple trivial zeros, positive nonvanishing
bounds at the edges of the critical strip, a small explicit zero-free
neighborhood of those edges, unbounded nontrivial zero counts, logarithmic
local multiplicity bounds, mesoscopic density, and effective convergence
of reciprocal-zero sums. Exact rational arithmetic certifies one simple
critical-line zero as a worked example.

The first-order statements concern specified rational approximants,
bounded searches and winding certificates. Complex analytic notation is
a readable interpretation of those statements. The Lean deliverable is
an explicit roadmap and object-theory interface, not a completed or
compiled Lean formalization. The reusable functional-equation draft is
included with its original proof-status qualifications.
\end{abstract}
\vfill
{\small Repository snapshot:
\code{VladimirReshetnikov/ProveIt} at\\
\code{b5a3ad037723c44d69a3d5b28a950c554327a989}.\\
Constants are deliberately coarse. No optimal zero-free region, critical-line
proportion, or Riemann hypothesis is asserted.}
\end{titlepage}
\clearpage\pagenumbering{roman}\hypersetup{pageanchor=true}
\begingroup\small\setlength{\parskip}{0pt}\tableofcontents\endgroup
\clearpage\pagenumbering{arabic}

\section{Results and the meaning of a constructive zero count}
\label{sec:results}

Write $s=\sigma+\ii t$ and use the normalization
\begin{equation}
 \xi(s)=\frac12s(s-1)\pi^{-s/2}\Gamma(s/2)\zeta(s).
 \label{eq:xi-definition}
\end{equation}
The completed function is entire, satisfies
$\xi(s)=\xi(1-s)$ and $\xi(\bar s)=\overline{\xi(s)}$, and has
$\xi(0)=\xi(1)=1/2$. Its zeros are exactly the nontrivial zeros of zeta,
with the same multiplicities. These standard normalizations agree with
\cite{DLMFreflection,DLMFzeros}. The functional equation is reused from
the quantitative human proof in \cite{RepoFE}; its complete archive is
included among the dependencies.

\subsection{The theorem inventory}
The following statements are established in their finite arithmetic
versions. Their usual complex-analytic readings are recorded here to make
the mathematical content immediately recognizable.

\begin{enumerate}
\item For rational $a>0$ and $\sigma\ge1+a$,
      $\abs{\zeta(\sigma+\ii t)}\ge a/(1+a)$.
      On the line $\sigma=1$ there are explicit positive lower bounds
      away from the pole. A uniform positive lower bound for
      $(s-1)\zeta(s)$ also covers a narrow neighborhood of that line.
\item The zeros in $\sigma\le0$ are exactly $-2,-4,-6,\ldots$.
      They are simple, and
      \[
       \zeta'(-2n)=
       \frac{(-1)^n(2n)!}{2(2\pi)^{2n}}\zeta(2n+1).
      \]
      All other zeros lie in $0<\sigma<1$; none is real.
\item Nontrivial zeros and their multiplicities are preserved by
      $s\mapsto\bar s$, $s\mapsto1-s$, and
      $s\mapsto1-\bar s$.
\item There is a fixed primitive recursive procedure which, on rational
      input $T\ge8$, returns a rational $U$ with $T<U<T+1$ and
      a certified zero-free boundary for
      $[-1,2]+\ii[0,U]$.
\item If $N(U)$ is the integer winding count on that boundary and
      \[
       \mathcal M(U)=\frac{U}{2\pi}\log\frac{U}{2\pi}-\frac{U}{2\pi},
      \]
      then
      \begin{equation}
       \boxed{\abs{N(U)-\mathcal M(U)}\le 8\log(U+4)+40.}
       \label{eq:main-rvm}
      \end{equation}
      The same estimate holds at every supplied regular height $U\ge8$.
\item The zero counts are unbounded with an explicit primitive recursive
      witness function. A bounded subdivision procedure also extracts
      a uniform sequence of pairwise apart primitive recursive zero names.
      The number of zeros in a unit-height strip and
      the multiplicity of any individual zero at height $T$ are
      $O(\log(T+4))$. Windows of length $\sqrt T$ have the expected
      asymptotic density $\log T/(2\pi)$.
\item For every integer $m\ge2$, the sum
      $\sum_\rho\rho^{-m}$ converges absolutely with an explicit
      primitive recursive tail bound. Pairing $\rho$ with $1-\rho$
      also makes the reciprocal sum converge absolutely in pairs.
\end{enumerate}
These are selected classical theorems about zero distribution; the
quantitative construction and arithmetic interpretation are the subject
of this report. Hardy's theorem on infinitely many critical-line zeros,
positive proportions on that line, sharp logarithmic zero-free regions,
and pair correlation are outside the theorem inventory.

\subsection{Why a count needs a certificate}
The predicate ``$\zeta(a+\ii b)=0$'' is not a decidable test on general
computable inputs. Nor may one silently search until a rational horizontal
line is found to avoid all zeros. The appropriate finite object is a
\emph{contour certificate}: rational domain vertices, rational
approximations to function values, a derivative bound, a strictly positive
separation margin, and the winding number of an associated rational
polygon. A checker recomputes the samples and verifies rational
inequalities.

The count means the total analytic multiplicity inside the contour.
It does not require choosing, ordering, or deciding equality of individual
zeros. Refinement produces certified clusters. In the ordinary complex
interpretation these counts equal the usual zero counts by the argument
principle, proved here through polynomial approximation.

Statements at arbitrary exact heights have a useful weaker representation.
For an integer $T\ge9$, construct regular heights
\[
 U_-\in(T-1,T),\qquad U_+\in(T,T+1).
\]
The two certified integers $N(U_-),N(U_+)$ bracket the count at $T$,
regardless of whether there is a zero exactly at $T$. Their difference is
bounded logarithmically. No convention for an undecidable boundary equality
is needed inside arithmetic.

\subsection{The first-order reading}
In this report $\ipra$ is first-order intuitionistic primitive recursive
arithmetic with all primitive recursive function symbols and
quantifier-free induction. Rational numbers, polynomials and certificates
are natural-number codes. Statements about fixed real or complex names
are schemas in specified primitive recursive definitions. They are not
quantification over a universe of arbitrary real numbers or over a
universal evaluator for all primitive recursive programs.

An estimate with error $\varepsilon$ means that every requested rational
precision admits the explicitly given approximants and rational error
budget. Precision and size parameters remain free during the finite
induction. The induction invariant is an equation, inequality, or Boolean
fold over bounded lists. Only after these inductions are complete are the
parameters universally generalized. Section~\ref{sec:foundation-audit}
identifies the exact repository proof type required by the Lean roadmap.

\section{Finite analytic representations and reusable identities}
\label{sec:analytic}

\subsection{Rational enclosures and elementary operations}
For $z=x+\ii y$ use
$\normone z=|x|+|y|$ and $|z|_\infty=\max(|x|,|y|)$.
These rational norms satisfy
\[
 |z|_\infty\le |z|\le\normone z\le2|z|,
 \qquad \normone{zw}\le\normone z\,\normone w.
\]
A complex ball is a rational centre $c$ and a rational error $r\ge0$,
with meaning $\normone{z-c}\le r$. For addition the radii add. For
multiplication an admissible radius is
\begin{equation}
 \normone c\,r_d+\normone d\,r_c+r_cr_d.
 \label{eq:ballproduct}
\end{equation}
Unless an estimate is explicitly in $\normone{\cdot}$, analytic
absolute-value bounds below use the complex Euclidean modulus. An
implementation that inserts one of them into an L1 radius first
multiplies it by two. The logarithm and exponential error estimates
used by the ball evaluator are direct L1 estimates.
If a centre is rounded to a dyadic rational, add its exact rational
rounding distance to the radius. Division is used only with a supplied
positive lower bound, or with an exactly nonzero rational denominator.

For rational $v>0$, reduce $v$ to $2^ku$, $1\le u<2$. The number of
reductions is bounded by the bit lengths of the numerator and denominator.
With $y=(u-1)/(u+1)$,
\begin{equation}
 \log u=2\sum_{j=0}^{L-1}\frac{y^{2j+1}}{2j+1}+r_L,
 \quad
 |r_L|\le\frac{2|y|^{2L+1}}{(2L+1)(1-y^2)}.
 \label{eq:log-series}
\end{equation}
Here $|y|\le1/3$. A linear bound on $L$ in the requested precision,
plus the bit length of $k$, therefore suffices. For a complex rational
$z$ of norm at most an integer $S\ge1$, the exponential series through
degree $d\ge2S$ has error
\begin{equation}
 \normone{\exp z-\sum_{j=0}^d z^j/j!}\le 2^{4S-d}.
 \label{eq:exp-series}
\end{equation}
Indeed the $2S$th term is at most $\exp S\le4^S$, and the subsequent
term ratios are at most $1/2$. Input perturbation of norm $r\le1$
costs at most $r4^{S+1}$. Both estimates are finite geometric-series
inequalities with their numerical parameters fixed.

The identity $\pi=4\arctan1$ with $2^{p+3}$ alternating terms is one
deliberately slow primitive recursive approximation to error $2^{-p}$.
Faster familiar formulas may be substituted after proving their error
contract. These choices make every displayed expression involving
$\pi$, real logarithms, exponentials, sine or cosine a specified
primitive recursive rational approximation.

\subsection{Euler summation and continuation}
Let $\{x\}=x-\lfloor x\rfloor$. Integrating on each interval
$[n,n+1]$ and summing gives, initially for $\sigma>1$,
\begin{equation}
 \zeta(s)=\frac{s}{s-1}
       -s\int_1^\infty\{x\}x^{-s-1}\,dx,
 \qquad \sigma>0,\quad s\ne1.
 \label{eq:euler-basic}
\end{equation}
The integral on the right has tail at most
$X^{-\sigma}/\sigma$. It consequently supplies a continuation with
an explicit convergence modulus on every region
$\sigma\ge a>0$, $|s|\le H$, $|s-1|\ge b>0$.
Differentiating $j$ times is justified by the bound
\begin{equation}
 \int_1^\infty(\log x)^j x^{-a-1}\,dx=\frac{j!}{a^{j+1}}.
 \label{eq:logmoments}
\end{equation}
One can avoid the word differentiation in the arithmetic proof:
expand $\exp(-h\log x)$ to the required finite degree, bound its
remainder, and integrate the resulting finite identities.

For wider disks and for the executable example we need the full finite
Euler--Maclaurin expression. Put
$(s)_j=s(s+1)\cdots(s+j-1)$ and let
$\widetilde B_j(x)=B_j(\{x\})$ denote the periodic Bernoulli polynomial.
For integers $N,m\ge1$ and $\sigma+2m>1$,
\begin{align}
 \zeta(s)={}&\sum_{n=1}^{N-1}n^{-s}
       +\frac{N^{1-s}}{s-1}+\frac12N^{-s}\notag\\
 &+\sum_{j=1}^m\frac{B_{2j}}{(2j)!}(s)_{2j-1}N^{1-s-2j}
       +R_{m,N}(s),\label{eq:em}\\
 R_{m,N}(s)={}&-\frac{(s)_{2m}}{(2m)!}
       \int_N^\infty\widetilde B_{2m}(x)x^{-s-2m}\,dx.
 \label{eq:em-remainder}
\end{align}
To prove this, integrate by parts on each unit interval, using
$B_j'=jB_{j-1}$ and $B_j(1)=B_j(0)$ for $j\ne1$. All inner endpoint
terms cancel. Repeat $2m$ times and then take a finite upper endpoint
to its explicitly bounded tail. This is a finite induction on the
number of integrations and on the number of unit cells.

The Bernoulli numbers are computed by
\[
 B_0=1,\qquad
 B_n=-\frac1{n+1}\sum_{j<n}\binom{n+1}{j}B_j.
\]
Thus
\begin{equation}
 C_m=\frac1{(2m)!}\sum_{j=0}^{2m}\binom{2m}{j}|B_j|
 \label{eq:bernoulli-bound}
\end{equation}
is a rational upper bound for
$\sup_{0\le x\le1}|B_{2m}(x)|/(2m)!$, obtained simply by bounding
each monomial. Consequently
\begin{equation}
 |R_{m,N}(s)|\le
 C_m\prod_{j<2m}(|\sigma+j|+|t|)
 \frac{N^{1-\sigma-2m}}{\sigma+2m-1}.
 \label{eq:em-bound}
\end{equation}
For the rational $\normone{\cdot}$-ball implementation, multiply the
right side by two. The rotation in $x^{-it}$ need not preserve the
$\normone{\cdot}$ norm, although it preserves the complex modulus.
For $\sigma\ge0$, replace the final factor by
$N^{1-2m}/(2m-1)$. This wholly rational bound is the one implemented in
\code{code/rational_zeta.py}. It requires no numerical library value of
zeta and no assumed Riemann hypothesis.

On a bounded rational box with $\sigma\ge-A$, choose $m=A+2$.
Then the exponent in the tail is at most $-A-3$. A power of two larger
than the rational coefficient bound times $2^{p+10}$ is already a
sufficient $N$ for error $2^{-p-10}$. Approximate the finite elementary
terms with a further budget for their rational coefficient magnitudes.
This gives a fixed primitive recursive approximation on every compact
box apart from $s=1$. To approximate the entire numerator everywhere,
multiply the finite formula symbolically by $s-1$ and cancel the
explicit pole term before evaluation. In particular, no evaluation of
an undefined zeta value at $s=1$ is performed.

\subsection{The functional-equation dependency}
For completeness we spell out the precise identities and quantitative
content reused from the included proof~\cite{RepoFE}. Define
\begin{align}
 \psi(x)&=\sum_{n\ge1}\exp(-\pi n^2x),\notag\\
 J(s)&=\int_1^\infty\psi(x)x^{s/2-1}\,dx,\notag\\
 \xi(s)&=\frac12+\frac{s(s-1)}2\bigl(J(s)+J(1-s)\bigr).
 \label{eq:theta-xi}
\end{align}
The theta relation
$1+2\psi(x)=x^{-1/2}(1+2\psi(1/x))$ follows there from a finite
Fej\'er calculation for the Gaussian, with Gaussian and Fourier tails
bounded before a limit is taken. Splitting the Mellin integral at $1$
gives~\eqref{eq:xi-definition} on $\sigma>1$ and its continuation.
Equation~\eqref{eq:theta-xi} gives both symmetries directly, with no
choice of a branch of the logarithm.

Here is a useful part of its explicit approximation schedule. On
$|\Re s|,|\Im s|\le B$, $B\ge2$, write
$\lambda(n)=\min\{j\le n:n\le2^j\}$ for $n\ge1$. Given $p$, put
\begin{align*}
 X&=p+\lambda(B!)+4,& K&=p+(B+1)\lambda(X)+4,\\
 D&=2(B+2)X^B,& q&=p+\lambda(XD)+4,&
 r&=p+\lambda(X)+4.
\end{align*}
Truncate the $x$ integral at $X$, the theta sum at $K$, use a left
mesh of width $2^{-q}$, and evaluate its finite summands so that the
error of their entire sum is at most $2^{-r}$ at each sample.
The four error bounds are
\begin{equation}
 2B!2^{-2X},\quad
 4X^{B+1}2^{-3(K+1)^2},\quad
 XD2^{-q},\quad X2^{-r}.
 \label{eq:theta-errors}
\end{equation}
Their sum is less than $2^{-p}$. These bounds follow respectively
from $\psi(x)\le2e^{-3x}$, a geometric estimate in the theta index,
a derivative bound for the integrand, and the sum of sample errors.
For example, $\psi\le1$ and
$|\psi'|\le4\sum_{n\ge1}n^2 8^{-n}=288/343<1$ justify the
displayed derivative budget. To obtain error $2^{-p}$ for $\xi$ itself,
first enclose both $s$ and $1-s$ in one common integer box of radius $B$,
and apply the integral schedule at precision
$p+\lambda(4B^2)+3$. This absorbs the two integral errors and the
multiplication by $s(s-1)/2$.
The same approximation parameters used for $s$ and $1-s$ make the
finite xi expression exactly symmetric.

The reciprocal gamma function is represented by
\begin{equation}
 G_N(z)=zN^{-z}\prod_{n=1}^N(1+z/n),\qquad G(z)=1/\Gamma(z).
 \label{eq:gamma-product}
\end{equation}
For integer $D\ge1$, $|z|\le D$, the included proof supplies
\begin{align}
 |G_N(z)|,|G(z)|&\le M_D:=D4^{D+D^2},\notag\\
 |G_N(z)-G(z)|&\le C_D/N,
 \quad C_D=2M_D(D+D^2),\quad N\ge2(D+D^2).
 \label{eq:gamma-modulus}
\end{align}
Away from the nonpositive integers it also gives a positive product
lower bound. The shift identity is $G(z)=zG(z+1)$.
These statements, their theta/Mellin identification, and the Gaussian
normalization $\Gamma(1/2)=\sqrt\pi$ are the reusable premises.
They are written proofs with finite numerical contracts, not completed
Lean theorems. Their source archive is preserved unchanged.

\subsection{From analytic formulas to rational polynomials}
Every use of a polynomial approximation below has a numerical error
requested in advance. On a closed disk strictly inside a domain of
holomorphy, a Taylor truncation is one construction: a bound $M_*$ on
a concentric radius $R_*$ gives
\[
 \sup_{|z-c|\le r}\left|f(z)-\sum_{j=0}^d a_j(z-c)^j\right|
 \le\frac{M_* (r/R_*)^{d+1}}{1-r/R_*}.
\]
The coefficients are obtained by finite Cauchy quadrature with an
explicit derivative bound. Alternatively one can use the fixed finite
Euler--Maclaurin or theta expressions above and replace each exponential
by its Taylor polynomial after approximating its real constants.
This second route avoids importing a general Taylor theorem: its degree,
coefficients and error budget follow by finite sum and product estimates.

All required disk sizes, larger-box bounds, degrees and working
precisions are fixed primitive recursive expressions in the input data.
The choice may be extremely inefficient. Efficiency is not used to
infer provability; the finite error estimates and induction invariants
are the reason this representation is suitable for $\ipra$.

\section{Zero exclusion with explicit positive margins}
\label{sec:elementary-zero-exclusion}

The distinction between a negative statement and a positive numerical
estimate matters here.  A proof of $\zeta(s)\ne0$ obtained by contradiction
does not, by itself, supply a primitive recursive precision at which a
numerical computation will certify nonvanishing.  We shall prove positive
rational lower bounds.  They are much weaker than the best analytic
bounds, but their role in the arithmetic construction is transparent.

We use the continuation formula
\begin{equation}
 I(s)=\int_1^\infty\{x\}x^{-s-1}\,dx,
 \qquad \zeta(s)=\frac{s}{s-1}-sI(s)
 \quad(\Re s>0,\ s\ne1).
 \label{eq:elem-euler-integral}
\end{equation}
Here and below an integral is shorthand for its finite quadrature
approximants and their stated tail estimates.  The braces in the
integrand cause no problem: split the integral at integer endpoints,
where its pieces are the smooth functions $(x-n)x^{-s-1}$.
For $\Re s\ge a>0$,
\begin{equation}
 |I(s)|\le a^{-1},\qquad |I'(s)|\le a^{-2},\qquad
 \left|\int_X^\infty\{x\}x^{-s-1}\,dx\right|
 \le \frac{X^{-a}}a .
 \label{eq:elem-I-bounds}
\end{equation}
The derivative estimate follows by inserting $\log x$ and using
$\int_1^\infty (\log x)x^{-a-1}dx=a^{-2}$; finite integration by parts
proves this last identity with its exponentially decreasing remainder.
Differentiating the finite approximants and their remainder bounds gives
\begin{equation}
 \zeta'(s)=-\frac1{(s-1)^2}-I(s)-sI'(s).
 \label{eq:elem-zeta-derivative}
\end{equation}

\subsection{The Euler region}

\begin{theorem}[Positive Euler-region nonvanishing]
\label{thm:elem-euler}
For $s=\sigma+it$ with $\sigma>1$,
\begin{equation}
 |\zeta(s)|\ge\frac{\sigma-1}{\sigma}.
 \label{eq:elem-euler-apartness}
\end{equation}
In particular, a supplied rational gap $\sigma\ge1+1/q$, $q\ge1$,
gives the rational lower bound $1/(q+1)$.
\end{theorem}

\begin{proof}
Define $D(s)=\sum_{n\ge1}\mu(n)n^{-s}$.  Trial division computes
$\mu(n)$ by primitive recursion.  Expanding the product over the distinct
prime divisors of $n$ gives the finite identity
$\sum_{d\mid n}\mu(d)=[n=1]$.  Consequently
\[
 \zeta(s)D(s)=1.
\]
For clarity, the passage from the finite convolution to this equality
has an explicit bound.  A product tail with $ab>N^2$ is contained in the
union $a>N$ or $b>N$.  Its absolute value is at most
\[
 2\zeta(\sigma)\frac{N^{1-\sigma}}{\sigma-1}.
\]
Thus the divisor calculation and this estimate prove the asserted
identity at every requested precision.  Since
\[
 |D(s)|\le\zeta(\sigma)
 \le1+\int_1^\infty x^{-\sigma}\,dx
 =\frac{\sigma}{\sigma-1},
\]
the result follows by multiplying the two norms.  No test for equality
of a complex number to zero is involved.
\end{proof}

\subsection{The finite 3--4--1 inequality}

\begin{lemma}[Euler positivity]
\label{lem:elem-341}
For $\sigma>1$ and real $t$,
\begin{equation}
 \zeta(\sigma)^3\,|\zeta(\sigma+it)|^4
                 |\zeta(\sigma+2it)|\ge1.
 \label{eq:elem-341}
\end{equation}
\end{lemma}

\begin{proof}
Put
\[
 L(s)=\sum_p\sum_{k\ge1}\frac{p^{-ks}}k.
\]
The sum restricted by $p^k\le N$ has error at most
$N^{1-\sigma}/(\sigma-1)$: each prime power is a distinct integer, and
its coefficient $1/k$ is at most one.  The power series for
$-\log(1-z)$, applied to finite Euler products, gives
$\exp L(s)=\zeta(s)$.  This equality also has an elementary tail proof:
expand each finite geometric product, and observe that the omitted
integers have a prime divisor beyond the included prime bound and hence
exceed that bound.

For every finite prime-power sum, use
\[
 3+4\cos u+\cos(2u)=2(1+\cos u)^2\ge0
\]
with $u=kt\log p$, multiplied by $p^{-k\sigma}/k$.  Summing and then
using the explicit tails gives
\[
 3L(\sigma)+4\Re L(\sigma+it)+\Re L(\sigma+2it)\ge0.
\]
Exponentiation proves \eqref{eq:elem-341}.  Equivalently, at precision
$2^{-m}$ the finite statement has lower bound $1-2^{-m}$ after the
error budget for exponentiation and multiplication is included.  For a
supplied gap $\sigma-1\ge1/q$, the choice
$N=2^{q(m+q)}$ already bounds the logarithm tail by $2^{-m}$, because
$q2^{-q}\le1$.  Enlarging $m$ by the explicitly computed multiplication
budget completes the finite version.
\end{proof}

\subsection{A uniform zero-free neighborhood of the one-line}

Let $F(s)=(s-1)\zeta(s)$, with the removable value $F(1)=1$.

\begin{theorem}[A rational uniform clearance bound]
\label{thm:elem-F-clearance}
For every integer $H\ge1$, put
\[
 \varepsilon_H=2^{-40}(H+1)^{-5},\qquad
 b_H=2^{-36}(H+1)^{-4}.
\]
Then
\begin{equation}
 |F(\sigma+it)|\ge b_H
 \quad\text{if}\quad |t|\le H,\quad
 |\sigma-1|\le\varepsilon_H.
 \label{eq:elem-F-clearance}
\end{equation}
In particular $\zeta$ has no zero on $\Re s=1$, and has no zero in
the displayed neighborhood of that line; its pole is represented by
the nonzero removable value of $F$.
\end{theorem}

\begin{proof}
First suppose $1/8\le|t|\le H$.  Set
\[
 M=128(H+1),\qquad B=16(H+1),\qquad
 h=(128BM^4)^{-1},\qquad\delta=Mh.
\]
For $1/2\le\sigma\le2$,
\eqref{eq:elem-I-bounds}--\eqref{eq:elem-zeta-derivative} give
\[
 |\zeta'(\sigma+it)|
 \le64+2+4(H+2)=4H+74\le M.
\]
For $1\le\sigma\le2$, the same integral formula gives
\[
 |\zeta(\sigma+2it)|
 \le(2+2H)(4+1)=10(H+1)\le B.
\]
Indeed $|\sigma+2it-1|\ge1/4$ and $|I(\sigma+2it)|\le1$.
Also $0<h\le1$ and
$\zeta(1+h)\le1+1/h\le2/h$.  Apply
Lemma~\ref{lem:elem-341} at $\sigma=1+h$ to obtain
\[
 |\zeta(1+h+it)|^4
 \ge\frac{h^3}{8B}=16M^4h^4=16\delta^4.
\]
The fourth power is increasing on nonnegative quantities, so
$|\zeta(1+h+it)|\ge2\delta$.  The derivative bound along the
horizontal segment now gives
\[
 |\zeta(1+it)|\ge2\delta-Mh=\delta.
\]
On the larger segment $|\sigma-1|\le h/2$ it gives
\begin{equation}
 |\zeta(\sigma+it)|\ge\delta-Mh/2=\delta/2.
 \label{eq:elem-zeta-near-one}
\end{equation}
The constants simplify to
\[
 \delta=2^{-32}(H+1)^{-4},\qquad h/2=\varepsilon_H.
\]

We next handle the pole and small ordinates directly.  From
\eqref{eq:elem-euler-integral},
\[
 F(s)=s-s(s-1)I(s).
\]
If $|s-1|\le r<1$, then
\[
 |F(s)-1|\le r+\frac{r(1+r)}{1-r}
 =\frac{2r}{1-r}.
\]
In particular
\begin{equation}
 |F(s)|\ge1/3\qquad (|s-1|\le1/4).
 \label{eq:elem-pole-disk}
\end{equation}
For $|t|\le3/16$ and $|\sigma-1|\le\varepsilon_H$, we have
$|s-1|\le1/4$, since $\varepsilon_H\le1/16$.  For
$|t|\ge1/8$, \eqref{eq:elem-zeta-near-one} instead gives
\[
 |F(s)|\ge\tfrac18\,\tfrac\delta2
 =\delta/16=b_H.
\]
These two alternatives overlap.  They can be selected constructively:
approximate $|t|$ with error at most $1/64$, and compare the rational
answer with $5/32$.  The smaller alternative implies $|t|<3/16$;
the larger implies $|t|>1/8$.  This is a bounded rational comparison,
not an invocation of trichotomy for arbitrary reals.  Both alternatives
give \eqref{eq:elem-F-clearance}.
\end{proof}

This theorem supplies an elementary polynomial-width zero-free region.
It does not assert the stronger classical width $c/\log(|t|+2)$.
The numerical constants were chosen to expose the arithmetic witnesses,
not to optimize the region.  On the simpler band $1\le|t|\le H$, the
same proof with $M=16(H+1)$ and $B=4(H+1)$ gives the better bounds
\[
 |\zeta(1+it)|\ge2^{-21}(H+1)^{-4},\qquad
 |\zeta(\sigma+it)|\ge2^{-22}(H+1)^{-4}
\]
whenever $|\sigma-1|\le2^{-26}(H+1)^{-5}$.

\section{Trivial zeros, the critical strip, and symmetry}
\label{sec:elementary-location}

We use the entire completion in its explicitly convergent form
\begin{align}
 \psi(u)&=\sum_{n\ge1}e^{-\pi n^2u},\notag\\
 J(s)&=\int_1^\infty\psi(u)
       \bigl(u^{s/2-1}+u^{-(s+1)/2}\bigr)\,du,\notag\\
 \xi(s)&=\frac12+\frac{s(s-1)}2J(s).
 \label{eq:elem-xi-integral}
\end{align}
The effective Gaussian identity and its Mellin calculation, established
in the analytic setup, identify this with
$\tfrac12s(s-1)\pi^{-s/2}\Gamma(s/2)\zeta(s)$ in $\Re s>1$.
The displayed integral itself gives
\begin{equation}
 \xi(s)=\xi(1-s),\qquad
 \xi(\bar s)=\overline{\xi(s)},\qquad
 \xi(0)=\xi(1)=1/2.
 \label{eq:elem-xi-symmetries}
\end{equation}
These identities hold between the explicit approximants with their
error bounds, so an abstract identity theorem is unnecessary.

\subsection{Reciprocal gamma factors and their zeros}

\begin{lemma}[Constructive reciprocal gamma zero classification]
\label{lem:elem-gamma}
The entire reciprocal gamma function
\begin{equation}
 G(z)=z e^{\gamma z}\prod_{n=1}^\infty(1+z/n)e^{-z/n}
 \label{eq:elem-gamma-product}
\end{equation}
has precisely the simple zeros $0,-1,-2,\ldots$, and
$G'(-n)=(-1)^n n!$.  On a rational bounded region with supplied positive
separations from these integers, the product supplies an explicit
positive lower bound for $|G|$.
\end{lemma}

\begin{proof}
For $|z|\le R$ and an integer $N>2R$, the logarithm of the tail of
\eqref{eq:elem-gamma-product} satisfies
\begin{equation}
 \left|\sum_{n>N}\bigl(\log(1+z/n)-z/n\bigr)\right|
 \le2R^2\sum_{n>N}n^{-2}\le2R^2/N.
 \label{eq:elem-gamma-tail}
\end{equation}
This uses only the power series for the logarithm in $|z/n|\le1/2$;
there is no choice of a global logarithm branch.  The tail is the
exponential of a number of bounded modulus and is therefore apart from
zero.  Every finite initial factor whose linear term is separated from
zero is apart from zero as well.

The classification does not infer a disjunction from a product being
zero.  Here is a bounded procedure which avoids that invalid general
constructive inference.  Approximate $z$ by a rational complex number
$z_0$ with error at most $1/100$.  Search the finite list
$0,-1,\ldots,-N$, where $N>2R$, for a center $-n$ with
$|z_0+n|\le1/4$, using squared distances.  If there is no such center,
every relevant factor has modulus at least $1/5$, after its explicit
integer denominator is included.  The finite product and
\eqref{eq:elem-gamma-tail} then certify $G(z)$ apart from zero.  If
there is such a center, it is unique, and $|z+n|<1/3$.  All other
linear factors are separated from zero by at least $2/3$ before their
integer denominators are included.  Thus
$G(z)=(z+n)H_n(z)$ on this disk, where $H_n$ has an explicit positive
lower bound.  Consequently $G(z)=0$ implies $z=-n$.

For example, off all the disks, the lower bound is a finite product of
$1/5$, the numbers $1/(5j)$ for $1\le j\le N$, and a rational lower
bound for $\exp(-R(N+1)-2R^2/N)$.  Such a lower bound follows from
$e<4$ by an integer power of $1/4$.  The disk case omits the one
distinguished factor and is otherwise identical.

For completeness, these normalizations are obtained directly from the
finite products.  Put $H_N=\sum_{j=1}^N1/j$.  At $z=1$ the finite
product is $e^\gamma(N+1)e^{-H_N}$, which tends to one with the
explicit harmonic-sum estimate
$|H_N-\log N-\gamma|\le1/N$.  If $G_N$ denotes the finite product,
cancellation after multiplication, as a polynomial identity, gives
\[
 zG_N(z+1)=(N+z+1)e^{\gamma-H_N}G_N(z).
\]
The factor on the right tends to one with a uniform error on every
supplied bounded region.  Passing to the explicit limits gives the
entire identity $G(z)=zG(z+1)$ without dividing by a possibly zero
factor.
Thus $G(1)=1$ and the shift identity have explicit finite proofs.
Iterating the shift through the finitely many factors gives
\[
 G(z)=z(z+1)\cdots(z+n)G(z+n+1).
\]
At $z=-n$, the derivative is $(-1)^n n!$.  Its positive magnitude
proves simplicity.
\end{proof}

Away from $s=0,1$ the continuation is equivalently
\begin{equation}
 \zeta(s)=\frac{2\pi^{s/2}\xi(s)G(s/2)}{s(s-1)}.
 \label{eq:elem-zeta-from-xi}
\end{equation}
This identity agrees with the right-half-plane definition and extends
it by the effective constructions.  In particular it is not an
additional nonvanishing assumption about gamma.

\subsection{Trivial zeros and strip localization}

\begin{theorem}[The trivial zeros are simple]
\label{thm:elem-trivial}
In $\Re s<0$, the zeros of $\zeta$ are exactly $s=-2n$, $n\ge1$.
They are simple, with
\begin{equation}
 \zeta'(-2n)=(-1)^n
 \frac{(2n)!}{2(2\pi)^{2n}}\zeta(2n+1),\qquad
 |\zeta'(-2n)|\ge\frac{(2n)!}{2\,8^{2n}}>0.
 \label{eq:elem-trivial-derivative}
\end{equation}
The disk $|s+2n|\le1/4$ contains no other zero.
\end{theorem}

\begin{proof}
For $\Re s<0$, the point $1-s$ lies in $\Re s>1$.
Theorem~\ref{thm:elem-euler}, the gamma bounds, and
\eqref{eq:elem-xi-symmetries} show that $\xi(s)=\xi(1-s)$ is apart
from zero.  Every other factor in \eqref{eq:elem-zeta-from-xi} except
$G(s/2)$ is apart from zero.  Lemma~\ref{lem:elem-gamma} therefore
gives precisely the listed zeros, and its local factorization gives
their simplicity and isolation on the stated disks.

For the derivative, differentiate just the vanishing factor in
\eqref{eq:elem-zeta-from-xi}.  Since $G'(-n)=(-1)^n n!$,
\[
 \zeta'(-2n)
 =\frac{(-1)^n n!\,\pi^{-n}\xi(1+2n)}{2n(2n+1)}.
\]
Substitute the right-half-plane expression for $\xi(1+2n)$ and use
\[
 \Gamma(n+1/2)=\frac{(2n)!}{4^n n!}\sqrt\pi.
\]
This last identity is obtained by $n$ applications of the gamma
recurrence to $\Gamma(1/2)=\sqrt\pi$; a general gamma reflection
formula is unnecessary.  The resulting expression is
\eqref{eq:elem-trivial-derivative}.  Finally $2\pi<8$ and
$\zeta(2n+1)\ge1$ give its rational positive lower bound.
\end{proof}

At $s=0$, the cancellation $G(s/2)/s\longrightarrow1/2$ in
\eqref{eq:elem-zeta-from-xi} gives $\zeta(0)=-1/2$.
Theorem~\ref{thm:elem-F-clearance} and reflection show that $\xi$ is
apart from zero on the imaginary axis as well as on the one-line.
The reciprocal gamma factor has no additional zero there.  These
observations exclude the boundary from the nontrivial zero set.

\begin{theorem}[Quantitative critical-strip localization]
\label{thm:elem-strip}
Define a nontrivial zero as a zero of the entire completion $\xi$.
If $\xi(s)=0$ and a supplied integer $H\ge1$ satisfies $|\Im s|\le H$,
then, with $\varepsilon_H=2^{-40}(H+1)^{-5}$,
\begin{equation}
 \frac{\varepsilon_H}{2}<\Re s
 <1-\frac{\varepsilon_H}{2}.
 \label{eq:elem-quantitative-strip}
\end{equation}
Inside this strip, the zeros and multiplicities of $\xi$ and $\zeta$
coincide.
\end{theorem}

\begin{proof}
In the band $|\Re s-1|\le\varepsilon_H$, write
\[
 \xi(s)=\tfrac12s\pi^{-s/2}\Gamma(s/2)F(s).
\]
The factors are apart from zero by the gamma product bounds and
Theorem~\ref{thm:elem-F-clearance}.  In the remaining right-hand
region $\Re s\ge1+\varepsilon_H/2$, the Euler lower bound applies,
again with apartness of every other factor.  The overlapping regions
thus give apartness whenever $\Re s\ge1-\varepsilon_H$.

To obtain the strict conclusion rather than only its double negation,
approximate $\Re s$ within $\varepsilon_H/8$ and compare the rational
answer with $1-3\varepsilon_H/4$.  The larger alternative places $s$
in the certified nonzero region and contradicts $\xi(s)=0$.  The
smaller alternative supplies $\Re s<1-\varepsilon_H/2$ directly.
Reflection $s\mapsto1-s$ gives the lower bound.  Finally
\eqref{eq:elem-zeta-from-xi} has an analytic invertible factor on
$0<\Re s<1$, proving the multiplicity assertion.
\end{proof}

\subsection{The real axis and an explicit initial zero-free rectangle}

\begin{lemma}[Positive real-axis completion]
\label{lem:elem-real-positive}
For every real $x$, $\xi(x)\ge3/8$.  For a positive integer $R$ and
$|s|\le R$,
\begin{equation}
 |\xi'(s)|\le D_R
 :=2(2R+1)R!+R(R+1)(R+1)!.
 \label{eq:elem-xi-derivative-bound}
\end{equation}
In particular $D_2=56$ and
\begin{equation}
 |\xi(s)|\ge1/4\quad
 (0\le\Re s\le1,\quad |\Im s|\le1/448).
 \label{eq:elem-initial-rectangle}
\end{equation}
\end{lemma}

\begin{proof}
For $u\ge1$, $\pi>3$ and $n^2\ge n$ give
\[
 0\le\psi(u)\le\sum_{n\ge1}e^{-3nu}
 =\frac{e^{-3u}}{1-e^{-3u}}\le2e^{-3u}.
\]
For real $x$ outside $(0,1)$, the correction in
\eqref{eq:elem-xi-integral} is nonnegative, so $\xi(x)\ge1/2$.
For $0\le x\le1$, each of the two powers of $u$ in $J(x)$ is at
most one.  Hence
\[
 0\le J(x)\le4\int_1^\infty e^{-3u}\,du
 =\tfrac43e^{-3}<1.
\]
Since $x(x-1)\ge-1/4$, we obtain $\xi(x)\ge1/2-1/8=3/8$.
These inequalities can first be checked for rational $x$, where the
case distinction is decidable, and then passed with their uniform
continuity estimates to a named real.  No decision about the sign of
an arbitrary real is required.

For $|s|\le R$, both complex powers in $J(s)$ have modulus at most
$u^R$.  Using $e^{-3u}\le e^{-u}$ and $\log u\le u$ gives
\[
 |J(s)|\le4\int_0^\infty u^Re^{-u}\,du=4R!,
 \qquad |J'(s)|\le2(R+1)!.
\]
Differentiate \eqref{eq:elem-xi-integral} and use
$|2s-1|\le2R+1$ and $|s(s-1)|\le R(R+1)$ to obtain
\eqref{eq:elem-xi-derivative-bound}.  The factorial integrals are
proved by a finite number of integrations by parts with an explicit
tail.

The vertical segment from $\sigma$ to $\sigma+it$, for
$0\le\sigma\le1$ and $|t|\le1/448$, stays inside $|s|\le2$.
The derivative bound therefore yields
\[
 |\xi(\sigma+it)|\ge\xi(\sigma)-56|t|
 \ge3/8-56/448=1/4.
\]
This proves \eqref{eq:elem-initial-rectangle}.
\end{proof}

\begin{corollary}[Symmetry, with multiplicity]
\label{cor:elem-symmetry}
If $\rho$ is a nontrivial zero of multiplicity $m$, then
$\bar\rho$, $1-\rho$, and $1-\bar\rho$ are zeros of the same
multiplicity.  There are no nontrivial real zeros.
\end{corollary}

\begin{proof}
The two identities in \eqref{eq:elem-xi-symmetries} carry the local
Taylor expansion at $\rho$ to those at the three listed points and
preserve the first nonzero coefficient.  Equivalently they carry the
local factorization and its positive leading-coefficient margin to the
reflected factorizations.  The real-axis assertion is
Lemma~\ref{lem:elem-real-positive}.  Some of the listed points may
coincide; the conclusion does not assume four distinct zeros.
\end{proof}

\subsection{Arithmetic content of the estimates}

For each supplied integer height bound, all constants above are
rational expressions in integer powers and factorials.  The Euler,
prime-power, and gamma products have explicit primitive recursive
truncation lengths.  The horizontal and vertical integral bounds use
finite quadrature estimates, not a choice of an intermediate point.
The few distinctions involving named reals are implemented by
overlapping rational comparisons.  Thus the statements have finite
arithmetical versions with no hidden unbounded search.  Their full
object-theory formalization still requires the derivations of the
finite analytic identities and error bounds in the chosen iPRA
calculus; an ambient Lean theorem about classical complex numbers is
not a substitute for those derivations.

\section{A uniform finite engine for rational complex polynomials}
\label{sec:polynomial-algebra}

The minimum-modulus argument requires factorization of a rational complex
polynomial whose degree varies with the requested accuracy. The repository's
FTA draft only asserts its full-factorization schema for each fixed degree.
That schema does not supply the uniform construction needed here. We give a
separate finite construction for rational polynomial codes. Its only
topological ingredient is the elementary winding calculus for rational
polygons: additivity under subdivision, invariance under a homotopy with a
positive separation margin, and winding one for a positively oriented square
about an interior point. These finite facts do not use the fundamental
theorem of algebra, Jensen's formula, or a general argument principle.

\subsection{Coefficient bounds and the squarefree reduction}

For a rational or Gaussian-rational polynomial $A(X)=\sum a_jX^j$, put
\[
 \mathcal B(A,B)=1+\sum_j |a_j|_1B^j,\qquad
 \mathcal B_r(A,B)=1+\sum_{j\ge r}
       j(j-1)\cdots(j-r+1)|a_j|_1B^{j-r},
\]
where $B\ge1$ is rational and $|x+iy|_1=|x|+|y|$.
These are rational upper bounds for the polynomial and its indicated
derivatives on $|z|\le B$. All ensuing constants can be computed exactly;
there is no appeal to an unspecified bound on coefficient heights.

The zero polynomial is rejected by the input guard; a nonzero constant
returns an empty root list and itself as leading coefficient. For a
nonconstant input first divide by its nonzero leading coefficient $a$,
so that the polynomial used in the gcd loop is monic.
Euclidean division and the extended Euclidean algorithm over $\mathbb Q(i)$
are bounded algorithms: a division has at most the current degree plus one
steps, and a nonzero Euclidean remainder strictly decreases the degree.
Leading coefficients are rational complex values, so their nonvanishing is
decidable and their inverses are exact rational operations. The usual gcd
identities give the squarefree decomposition
\[
                  P=a\prod_{j=1}^{d}S_j^{\,j},
\]
with monic, squarefree, pairwise coprime $S_j$, allowing $S_j=1$.
For clarity, writing $P_0=P/a$, this can be computed by setting
$c=\gcd(P_0,P_0')$, $w=P_0/c$ and,
successively, $y=\gcd(w,c)$, $S_j=w/y$, $w=y$, $c=c/y$.
The loop stops when $w=1$ and fills remaining entries with $1$; a zero
Euclidean remainder stops that particular gcd computation. The outer
loop runs at most $d=\deg P$ times. Polynomial identities and the
degree decrease prove the decomposition; its implementation can verify the
displayed product directly. Thus it suffices to treat a monic squarefree
polynomial $S$ of degree $d\ge1$.

Compute polynomials $U,V$ with
\begin{equation}
                       US+VS'=1.                    \label{eq:poly-bezout}
\end{equation}
Let $R$ be an integer greater than $2+\sum_{j<d}|s_j|_1$, and set
$B=4(R+1)$. Leading-term dominance gives $|S(z)|\ge1$ whenever
$|z|\ge R-1$. For a root name $\alpha$, read a rational approximation
$a$ with error at most $1/4$ and decide the rational comparison
$|a|\le R-1/2$ by squaring. The opposite branch would imply
$|\alpha|>R-3/4>R-1$ and contradict the displayed lower bound.
Consequently $|\alpha|\le R-1/4<R$. This obtains a strict bound with
a margin, rather than inferring it from the negation of an order relation.
Put
\[
 A=\mathcal B(U,B),\quad K=\mathcal B(V,B),\quad
 M_1=\mathcal B_1(S,B),\quad M_2=\mathcal B_2(S,B).
\]
In particular $A,K,M_1,M_2\ge1$. At a root $\alpha$,
\begin{equation}
                  |S'(\alpha)|\ge K^{-1}.            \label{eq:poly-root-derivative}
\end{equation}
The finite Taylor identity gives
\[
 |S(\beta)-S(\alpha)-S'(\alpha)(\beta-\alpha)|
       \le\tfrac12 M_2|\beta-\alpha|^2.
\]
Consequently two roots at distance at most $(KM_2)^{-1}$ must be equal:
the last inequality and \eqref{eq:poly-root-derivative} imply
$|\beta-\alpha|/K\le |\beta-\alpha|/(2K)$.
This proves the needed separation for roots known to be apart without
dividing by an undecided difference of named numbers.

\subsection{Real roots by a bounded grid and bisection}

The following real procedure will be used several times. Given a nonzero
rational polynomial $F$, remove repeated factors by gcd, normalize to monic,
and call the result $F_0$. Compute $U_0F_0+V_0F_0'=1$, a Cauchy root bound
$R_0$, and bounds $K_0\ge |V_0|$, $M_0\ge\max(1,|F_0''|)$ on
$[-R_0-1,R_0+1]$. Choose
\[
                       h=\frac1{64K_0M_0}.
\]
For grid centres $jh$ covering $[-R_0-h,R_0+h]$, inspect the overlapping
intervals $[jh-2h,jh+2h]$. Retain an interval precisely when its two rational
endpoint values of $F_0$ have strictly opposite signs. This is a finite
scan with fewer than $2R_0/h+10$ centres.

Every retained interval contains a uniquely specified real root name.
Bisection uses exact rational signs; an exactly zero midpoint produces a
constant name. After $q$ bisections the width is $4h2^{-q}$, so the number
of bisections for any requested accuracy is explicitly bounded. The nested
midpoints are Cauchy, and the polynomial Lipschitz bound shows their
residual tends to zero at the same explicit rate. Existence here is this
construction, rather than an invocation of a general intermediate-value
principle.

Conversely, every real root $a$ of $F_0$ lies in one of the retained
intervals. Choose a grid centre within $h$ of an approximation to $a$,
with the approximation fine enough to keep the distance below $3h/2$.
On that interval the derivative differs from $F_0'(a)$ by less than
$1/(4K_0)$, while $|F_0'(a)|\ge1/K_0$. The finite Taylor identity therefore
makes the endpoint signs strictly opposite. No exact comparison of $a$
with a mesh point is required. Distinct roots are at least
$(K_0M_0)^{-1}$ apart; any two overlapping retained intervals consequently
name the same root. Merge overlapping intervals and keep one representative
per group. This gives a finite list of pairwise apart root codes containing
every real root. Each code consists of $F_0$ and one strict sign-changing
rational interval. Its evaluator is the fixed bisection program just
described, uniform in the numerical polynomial code and the precision.

The identities $F_0\mid F$ and $F\mid F_0^{\deg F}$, obtainable during
the gcd calculation, prove that the list also describes precisely the
real zeros of the original $F$.

\subsection{Rational coordinate polynomials for complex roots}

Set $Q=S\overline S\in\mathbb Q[X]$, where conjugation acts on the
coefficients, and let $n=2d$. Consider the finite rational algebra represented
by arrays in the basis $u^av^b$, $0\le a,b<n$, with multiplication reduced
by the monic relations $Q(u)=Q(v)=0$. Its array dimension is $n^2$.
Form the multiplication matrices for
\[
                        x_0=(u+v)/2,
             \qquad     y_0^2=-(u-v)^2/4.
\]
Let $A_X$ be the characteristic polynomial of the first matrix and let
$A_Y(t)$ be the characteristic polynomial of the second matrix evaluated
at $t^2$. Both are nonzero monic rational polynomials, of degrees $n^2$
and $2n^2$, respectively. Determinants can be computed by the finite
permutation formula; efficiency is unnecessary here.

If $\alpha=x+iy$ is a root of $S$, then $Q(\alpha)=Q(\overline\alpha)=0$.
Evaluating the array algebra at $(u,v)=(\alpha,\overline\alpha)$ and using
Cayley--Hamilton gives
\begin{equation}
                         A_X(x)=A_Y(y)=0.              \label{eq:coordinate-polynomials}
\end{equation}
The use of Cayley--Hamilton is purely finite: the adjugate identity for
$tI-M$ gives its characteristic-polynomial identity, and reduction by the
two monic polynomials gives an identity of coefficient arrays. No
existence of complex roots has been used to construct these matrices.

Apply the real-root procedure to $A_X,A_Y$. Every possible complex root
of $S$ is now among a finite Cartesian product of coded real algebraic
numbers. The remaining task is an exact, bounded zero test for each pair.

\subsection{Exact zero and order tests on these codes}

Suppose $x,y$ are real root codes for monic rational polynomials $F,G$,
and $W(X,Y)$ is a rational polynomial. In the finite array algebra
\[
                  \mathbb Q[X,Y]/(F(X),G(Y)),
\]
compute the characteristic polynomial $C$ of multiplication by $W$.
It is monic and satisfies $C(W(x,y))=0$. Factor out its initial power of
the variable by examining its rational coefficients:
\[
                    C(t)=t^mC_0(t),\qquad C_0(0)\ne0.
\]
If $C_0$ is constant, $W(x,y)=0$. Otherwise write $C_0(t)=\sum c_jt^j$
and choose the positive rational
\begin{equation}
 \eta=\min\left(1,\frac{|c_0|}{2(1+\sum_{j\ge1}|c_j|)}\right).
                                                        \label{eq:algebraic-gap}
\end{equation}
Every nonzero real root of $C$ has absolute value at least $\eta$.
Indeed, for $|t|<\eta$, the finite geometric estimate gives
$|C_0(t)|\ge |c_0|/2$; its reciprocal is then permitted, and
$C(t)=0$ implies $t^m=0$ and hence $t=0$.

Evaluate $W(x,y)$ with error at most $\eta/8$, using the bisection codes
and the coefficient bound for the two-variable polynomial on their known
input boxes. This fixes a numerical precision before any evaluation loop.
For the rational result $w$, the test $|w|\le\eta/2$ means the exact
value is zero; otherwise the value is apart from zero and has the sign of
$w$. The gap \eqref{eq:algebraic-gap} proves both branches. The zero
branch has a proof at every precision by the displayed annihilating
identity, rather than by a real-equality oracle.

Apply the test to the real and imaginary parts of $S(x+iy)$. Retain
exactly the pairs passing both zero tests. Equations
\eqref{eq:coordinate-polynomials} and completeness of the real-root lists
prove that this is a complete list of actual roots. Different retained
pairs are apart because their real-root intervals distinguish at least
one coordinate. The same test, applied to $x^2+y^2-r^2$, decides whether
a root is inside, on, or outside a rational circle. Thus a later selection
of roots with $|\alpha|\le3$ includes boundary roots correctly.

\subsection{Why the complete list has exactly the degree}

The preceding construction lists all roots but has not assumed that the
list is nonempty. The following quantitative argument proves its length
is $d$, and closes that possible circularity.

We first need a bounded local root construction. Let $z$ be Gaussian
rational with $|z|\le B-1$, and suppose $|S(z)|\le\varepsilon$, where
\begin{equation}
 \varepsilon\le\min\left(
       \frac1{4A},\frac1{8K},\frac1{16K^2M_2},
       \frac b{16K}\right),\qquad 0<b\le1.              \label{eq:newton-budget}
\end{equation}
By \eqref{eq:poly-bezout}, $|S'(z)|\ge1/(2K)$. Set
\[
        r=4K\varepsilon,\qquad
        T(w)=w-S(w)/S'(z).
\]
The denominator is a nonzero Gaussian rational and is computed exactly.
On $|w-z|\le r$, the finite difference-of-powers formula gives
\[
 |T(w)-T(w')|\le2KM_2r\,|w-w'|\le\tfrac12|w-w'|,
       \qquad |T(z)-z|\le r/2.
\]
Thus the rational iteration $z,T(z),T^2(z),\ldots$ remains in that
disk and has a geometric Cauchy modulus. It gives a root name within
$r\le b/4$ of $z$. Its step count is the requested precision plus a
computed binary length; this is one primitive recursion, not an unbounded
Newton search. The coefficient identity also gives a geometric residual
bound for $S$ at the iterates.

Now put $r_0=(64KM_2)^{-1}$. Around every root in the complete list choose
a rational square of half-side $r_0$, whose centre approximates that root
with error at most $r_0/16$. These squares are disjoint, by root separation.
Their boundaries have distance at least $15r_0/16$ from their corresponding
roots. The Taylor inequality and \eqref{eq:poly-root-derivative} give a
zero-free homotopy of $S$ on each boundary to
$S'(\alpha)(z-\alpha)$: the remainder is less than $1/32$ of the linear
term. The image therefore has winding one.

Take $b=r_0/2$ in \eqref{eq:newton-budget}. At every rational point outside
the interiors of all these squares and inside the outer square of
half-side $R+1$, one has
\begin{equation}
                              |S(z)|>\varepsilon.       \label{eq:poly-outside-lower}
\end{equation}
Otherwise the just-proved rational iteration would construct a root within
$b/4$ of $z$. Completeness of the algebraic list identifies it with a listed
root, whose distance from every such $z$ is greater than $b$, a contradiction.
This conclusion uses a decidable rational comparison: the branch
$|S(z)|\le\varepsilon$ is eliminated by the explicit construction.

Partition the outer square minus the root squares by a rational rectangular
grid, including every square edge and further subdividing until every cell
has diameter at most $\varepsilon/(4M_1)$. All cutoffs are explicit. The
images of each cell's vertices, and the polynomial images of its edges,
stay in a disk of radius at most $\varepsilon/4$ about the image of one
vertex. By \eqref{eq:poly-outside-lower} that disk misses zero. Each cell
image has winding zero. Cancelling shared edges leaves
\[
 \operatorname{wind}(S\circ\partial Q_{R+1})
       =\sum_{\text{root squares}}1.
\]
On the outer square, the leading term of the monic $S$ dominates the sum
of the lower terms. The zero-free straight homotopy to $z^d$ gives winding
$d$. This proves that the list contains exactly $d$ roots.

Synthetic division now yields
\[
                             S(z)=\prod_{j=1}^d(z-\alpha_j).
\]
When a root is removed, every remaining root is still a root of the
quotient, since its difference from the removed root has an explicit
positive lower bound. Repeating the argument for the rational squarefree
factors $S_j$ gives the full multiset factorization of the original $P$,
including multiplicities.

\subsection{Primitive recursion and the induction audit}

Every output root is evaluated by the same fixed bounded bisection routine
on its finite polynomial-and-interval code. Full factorization therefore
does not iterate a higher-order root-selection operation of variable depth.
The operations are rational arithmetic, coefficient-array reduction,
degree-bounded Euclidean division, determinant expansion over a finite
permutation list, bounded grid scans, bisection with a specified number of
steps, and a geometrically bounded rational iteration. Their loop lengths
are obtained by composition, primitive recursion, factorial, powers,
rational ceilings, and bounded binary logarithms. Computed coefficient
bounds replace any hidden root-separation or height oracle.

For the internal proof, requested precisions are fixed before induction.
The invariants are the exact division identities, degree decreases,
matrix-array identities, rational brackets and polynomial signs, the
rational contraction and residual inequalities, and the finite signed-edge
balance. Bounded lists of such conditions are encoded by Boolean folds.
Synthetic-division identities and products of algebraic root names are
checked at an allocated finite precision, using root bounds and the finite
telescoping product inequality; all root evaluations use their uniform
finite codes. These are quantifier-free induction arguments.

This establishes the uniform finite polynomial engine required below.
It does not postulate a total evaluator for arbitrary primitive recursive
programs, a compactness principle, a general real zero test, or the
fundamental theorem of algebra as an unproved analytic premise.


\section{Finite winding, stability, and analytic zero counts}
\label{sec:finite-winding}

The polynomial algebra just established is used as a finite intermediate
calculus. It is not a claim that general computable complex numbers have
decidable equality. In particular, the algorithm below never asks whether
an unrepresented analytic zero is exactly on a boundary.

\subsection{Rational polygons and their winding}
Let $w_0,\ldots,w_n=w_0$ be a rational polygon which avoids the origin,
with explicit separation bounds on its edges. Its integer winding number
can be computed without logarithms. Choose a rational ray from zero that
contains no vertex and is parallel to no edge. Among more than twice the
number of vertices and edges, distinct rational slopes, at least one
avoids every forbidden direction. Test each by rational determinants.
Count oriented crossings of that ray, using the sign of a determinant
and a rational intersection parameter. Upward and downward crossings
contribute opposite signs; the orientation is fixed so a positively
oriented square about zero has winding one.

The elementary invariance rules have finite proofs. Inserting a vertex
on an edge preserves the crossing sum. Reversing an edge negates its
contribution. If a small triangular boundary stays in an open half-plane
through the origin, its sum is zero. Triangulate a supplied homotopy
whose images stay a positive distance from zero, and refine until each
small triangle has that property. Interior edges cancel. The derivative
bound and positive margin give the number of subdivisions before the
calculation starts. Thus homotopy invariance is an identity of finite
integer sums, with a quantitative refinement bound.

For two nonzero polygonal paths, winding is additive under pointwise
multiplication. This follows by subdividing until each image lies in a
disk of radius less than half its centre's distance from zero. On each
such disk the logarithm of $1+z$ is its explicit convergent series, and
argument increments add. Equivalently one can verify the corresponding
crossing identity on a finite triangulation. The same argument proves
additivity under concatenation and subdivision of a domain polygon.

\subsection{Passing from a fixed analytic function to a polygon}
Let a rational domain polygon $\gamma$ have segment length at most $h$,
and suppose the fixed function $f$ has derivative bounded by $B$ on
each segment. Suppose also $|f|\ge\mu>0$ on the entire boundary.
At its vertices compute rational values $w_j$ with
$|w_j-f(\gamma_j)|\le e$. Linear interpolation $W$ between the $w_j$
has error at most
\begin{equation}
 |W(t)-f(\gamma(t))|\le Bh/2+e.
 \label{eq:winding-interpolation}
\end{equation}
Indeed at a relative position $a\in[0,1]$ the two endpoint errors
contribute $e$, and the Lipschitz errors contribute
$2a(1-a)Bh\le Bh/2$. If $Bh/2+e<\mu$, the straight homotopy avoids
zero. Define the certified count to be $\wind(W,0)$.

This number is independent of any further certified refinement. For two
certificates use a common refinement and their error budgets to construct
the zero-free homotopy between their polygons. A certificate can also
verify a margin from its samples: if
$|w_j|\ge\mu+e+Bh$, every point on the corresponding edge has modulus
at least $\mu$. Squared Euclidean norms of rational samples are
decidable; every use of a non-rational norm in this argument is backed
by a rational upper or lower bound.

\subsection{The polynomial argument principle and Rouch\'e stability}
For a rational complex polynomial $P$, factorization in the finite
algebraic calculus gives
$P(z)=a\prod_j(z-\alpha_j)$. If a rational rectangle avoids the
algebraic roots, the multiplication rule for winding gives
\[
 \wind(P\circ\partial R,0)
 =\sum_j\wind(\partial R-\alpha_j,0).
\]
A translated rectangle has winding one around an interior point and
zero around an exterior point. These alternatives are decidable for
the algebraic roots and are proved by the crossing definition. Hence
the winding equals the number of roots in $R$, with multiplicities.
In particular it is nonnegative, is additive under domain subdivision,
and cannot exceed the degree.

Now approximate the fixed analytic $f$ by a rational polynomial $P$
on a neighborhood of a certified boundary, with error less than half
its clearance. The homotopy $(1-t)f+tP$ avoids zero on that boundary,
so its count is the polynomial count. The same holds simultaneously
for any finite family of certified subregions by taking the minimum
of their finitely many rational margins. This proves additivity,
nonnegativity and all counting comparisons used below.

In the usual analytic interpretation, this is the argument principle.
For an isolated zero with a supplied local factorization
$f(z)=(z-\rho)^m g(z)$ and an apartness bound for $g$, a small
surrounding contour has count $m$. Conversely positive count is the
constructive assertion that the region contains a zero cluster. Finer
certified grids localize such clusters. No use of the later counting
formula depends on selecting an exact zero or its exact multiplicity.
Section~\ref{sec:zero-extraction} gives a uniform point-name extraction
from positive count. The explicit simple-zero example in
Section~\ref{sec:example} also supplies its own direct contraction,
independent of that general construction.

Rouch\'e's theorem in the form used here is just the same homotopy
estimate. If $|g-f|<|f|$ on a boundary with a supplied positive
difference margin, $f$ and $g$ have the same integer winding. Quantify
the inequality first, choose the rational mesh second, and prove the
finite homotopy third. No choice of a ``sufficiently fine'' mesh is
left without a bound.

\subsection{Induction discipline}
The algebraic factorization, vertex enumeration, ray selection, crossing
sum, polynomial approximation and refinement are bounded programs.
For fixed requested precision, their correctness invariants are finite
polynomial identities, rational inequalities, or equations for integer
partial sums. A bounded list of inequalities is represented by its
Boolean conjunction. These are exactly the quantifier-free invariants
needed by $\ipra$. Analytic notation is eliminated by the schedules of
Section~\ref{sec:analytic} before induction is applied.

\section{The complex gamma estimate needed for counting}
\label{sec:stirling}

The real factorial Stirling estimate alone is insufficient for the
argument of $\Gamma(1/4+\ii T/2)$. We derive the needed complex estimate
directly. Let $\Log z$ be the branch with argument in
$(-\pi/2,\pi/2)$ on $\Re z>0$, and let $\log\Gamma z$ denote its
analytic logarithm which is real on the positive real axis. It can be
defined by the logarithm of the finite products in
\eqref{eq:gamma-product}; all factors have positive real part after
the initial $z$ factor.

\begin{lemma}[Complex Euler--Maclaurin remainder]
For $\Re z>0$,
\begin{equation}
 \log\Gamma z=(z-\tfrac12)\Log z-z+C+R(z),\qquad C\in\R,
 \label{eq:complex-stirling}
\end{equation}
where
\begin{equation}
 R(z)=\frac1{12z}-\frac12\int_0^\infty
          \frac{B_2(\{x\})}{(z+x)^2}\,dx,
 \qquad |R(z)|<\frac1{4|z|}.
 \label{eq:gamma-remainder}
\end{equation}
The constant is $C=\tfrac12\log(2\pi)$, but its value is not required
in this argument.
\end{lemma}

\begin{proof}
Apply two integrations by parts on each unit interval to
$f(x)=\Log(z+x)$. Since $f'=1/(z+x)$ and $f''=-1/(z+x)^2$, the finite
Euler--Maclaurin identity on $[0,n]$ is
\begin{align*}
 \sum_{j=0}^{n-1}\Log(z+j)
 ={}&\int_0^n\Log(z+x)\,dx
     +\frac{\Log z-\Log(z+n)}2\\
 &+\frac1{12}\left(\frac1{z+n}-\frac1z\right)
   +\frac12\int_0^n\frac{B_2(\{x\})}{(z+x)^2}\,dx.
\end{align*}
Insert this equality into
\[
 \log\Gamma z=
 \lim_{n\to\infty}\left(\log((n-1)!)+z\log n
                 -\sum_{j=0}^{n-1}\Log(z+j)\right).
\]
The elementary antiderivative of $\Log(z+x)$ yields
\eqref{eq:complex-stirling}--\eqref{eq:gamma-remainder}; the
remaining constant is the real limit of
$\log((n-1)!)-(n-1/2)\log n+n$. Existence with an $O(1/n)$ modulus
follows by the same finite formula at $z=1$. No complex limiting
constant has to be selected. The upper-endpoint errors are bounded
uniformly on $|z|\le D$, $\Re z\ge a>0$ by a rational multiple of
$(D+1)^2/n$, using the finite series for $\Log(1+z/n)$ when $n\ge2D$.

Since $|B_2(x)|\le1/6$ on $[0,1]$, it remains to bound a real
integral. Write $z=|z|e^{\ii\theta}$, $|\theta|<\pi/2$. Direct
integration gives, including the limiting case $\theta=0$,
\[
 \int_0^\infty\frac{dx}{|z+x|^2}
 =\frac{|\theta|}{|z|\sin|\theta|}
 \le\frac{\pi}{2|z|}.
\]
Therefore
\[
 |R(z)|\le\frac1{12|z|}+\frac{\pi}{24|z|}
       =\frac{2+\pi}{24|z|}<\frac1{4|z|}.
\]
All divisions are bounded away from zero on the stated half-plane
boxes. For the integral tail beyond $X$ one may use
$1/[12(X+a)]$. On each unit interval, the derivative of its integrand
is bounded by a rational expression in $a^{-1}$ and $|z|$; splitting
at the integer endpoints and taking a fixed rational mesh bounds
quadrature error. Thus the displayed proof has explicit primitive
recursive cutoff and mesh choices. It does not require an unspecified
asymptotic expansion.
\end{proof}

Define the continuous gamma argument
\begin{equation}
 \vartheta(T)=\Im\log\Gamma(\tfrac14+\tfrac{\ii T}2)
               -\frac T2\log\pi,
 \label{eq:theta-definition}
\end{equation}
with the branch just fixed. In particular $\vartheta(0)=0$.
The distinction between this continuous logarithm and the principal
argument of the gamma value matters: reducing modulo $2\pi$ would
destroy a zero-counting formula. The normalization agrees with the
standard Riemann--Siegel theta function; see
\cite{DLMFgamma,DLMFzeros}.

\begin{lemma}[Explicit theta estimate]
For $T\ge1$,
\begin{equation}
 \left|\vartheta(T)-
 \left(\frac T2\log\frac{T}{2\pi}-\frac T2-\frac\pi8\right)\right|
 <\frac1T.
 \label{eq:theta-bound}
\end{equation}
\end{lemma}

\begin{proof}
Put $a=1/4$, $b=T/2$, $r=(a^2+b^2)^{1/2}$ and
$\theta=\arg(a+\ii b)$. The imaginary part of the main expression
in~\eqref{eq:complex-stirling} is
\[
 b\log r-\tfrac14\theta-b.
\]
The replacement of $\log r$ by $\log b$ costs at most
\[
 \frac b2\log(1+a^2/b^2)\le\frac{a^2}{2b}=\frac1{16T}.
\]
The replacement of $\theta$ by $\pi/2$ costs at most
$\tfrac14\arctan(a/b)\le1/(8T)$. Finally
$|R(a+\ii b)|<1/(4b)=1/(2T)$. The total is less than
$11/(16T)<1/T$, and the real constant $C$ contributes no imaginary
part.
\end{proof}

% Standalone section fragment. Uses amsmath, amsthm, amssymb only.
\section{Bounded selection of regular heights}
\label{sec:regular-heights}

A sentence such as ``move the horizontal edge a little until it misses all
zeros'' is not an algorithm with a primitive recursive bound.  This section
supplies that bound.  All the constants are intentionally generous.
The polynomial facts used below concern rational complex polynomials; their
algebraic roots are represented by finite real-algebraic data.  In particular,
root multiplicities and comparisons with a rational circle are decidable by
bounded algebraic algorithms.  The uniform, variable-degree version of this
finite polynomial lemma is required here, not merely a separate theorem for
each fixed degree.

\subsection{A finite polynomial substitute for Jensen's inequality}

Let $P$ be a polynomial, $P(0)\ne0$, and suppose $|P(z)|\le B$ on
$|z|=R$.  If $a_1,\ldots,a_q$ are any of its roots in $|z|\le r<R$,
repeated according to multiplicity, put
\[
 B_a(z)=\frac{R(z-a)}{R^2-\overline a z},\qquad
 H(z)=P(z)\prod_{j=1}^q B_{a_j}(z)^{-1}.
\]
The apparent divisions cancel polynomial factors, so $H$ is a polynomial.
On $|z|=R$ each $|B_a(z)|=1$.  The constant-coefficient mean identity for a
polynomial therefore gives
\[
 |P(0)|\prod_{j=1}^q\frac R{|a_j|}=|H(0)|\le B,
 \qquad
 |P(0)|(R/r)^q\le B.                         \tag{J}
\]
There is no assumption that the outer circle avoids roots.  The mean identity
is verified by integrating the finitely many monomials, or by a sufficiently
fine finite Fourier sum.  Thus (J), unlike a qualitative invocation of
Jensen's formula, is already a finite algebraic inequality with explicit
approximation margins.

\subsection{A minimum-modulus estimate independent of the degree}

Suppose $P\in\mathbb Q(i)[z]$ satisfies
\[
 |P(0)|\ge\tfrac18,\qquad |P(z)|\le2M\quad(|z|\le4),\qquad M\ge1.
\]
Let $k$ be the least integer for which $2^k\ge16M$, and put $Q=4k$.
By (J), the number $q$ of roots in $|z|\le3$ is at most $Q$:
indeed $(4/3)^q\le16M\le2^k$, whereas $(4/3)^4>2$.
Remove exactly these roots with the radius-four factors $B_a$ above.
The resulting polynomial $H$ has no root in $|z|<3$, satisfies
$|H|\le2M$ on $|z|\le4$, and $|H(0)|\ge1/8$.

For completeness, the only harmonic estimate needed is the following
explicit version of Harnack's inequality.  On $|z|\le11/4$ the function
\[
 u(z)=\log\frac{2M}{|H(z)|}
\]
is nonnegative and harmonic.  The Poisson kernel at $|z|\le5/2$, relative to
the circle of radius $11/4$, lies between positive bounds and is at most
\[
 \frac{11/4+5/2}{11/4-5/2}=21.
\]
Consequently $u(z)\le21u(0)$ and
\[
 |H(z)|\ge \frac{1}{2^{83}M^{20}}\quad(|z|\le5/2).       \tag{H}
\]
This use of Poisson integration has a finite implementation.  If $d=\deg H$,
all roots of $H$ have modulus greater than three, whence its factorization
already gives
\[
 |H(z)|\ge |H(0)|\,12^{-d}\quad (|z|\le11/4).
\]
That inferior but explicit lower bound controls the quadrature and logarithm
errors while proving the much stronger bound (H).  More concretely, write
$H(z)=H(0)\prod_{\nu=1}^d(1-z/b_\nu)$, where $|b_\nu|>3$.
On the radius-$11/4$ circle the logarithm series
$\log(1-z/b_\nu)=-\sum_{j\ge1}(z/b_\nu)^j/j$ has ratio at most $11/12$.
Truncation after $K$ terms has total error at most
$12d(11/12)^{K+1}$, and the angular derivative of $\log|H|$ is at most
$11d$.  Finite Fourier identities prove the mean and Poisson identities
for the truncated sums; the displayed geometric tail and the positive
lower bound permit passage to any requested rational error.
Their dependence on $d$ is harmless: the approximation degree has already
been chosen.  They do not enter the approximation tolerance used to choose
that degree.

If $|z-a_j|\ge\delta$ for every removed root and $|z|\le5/2$, then
\[
 |B_{a_j}(z)|\ge\frac{4\delta}{16+3(5/2)}>\frac\delta6.
\]
For $0<\delta\le1$ this proves
\[
 |P(z)|\ge \frac{1}{2^{83}M^{20}}(\delta/6)^Q
                                                        \tag{MM}
\]
outside at most $Q$ disks of radius $\delta$.  The exponent is controlled by
$\log M$, not by the approximation degree.  This distinction eliminates
an otherwise circular choice of precision.

\subsection{The regular-height algorithm}

Take rational $T\ge8$, write $A=\lceil T\rceil$, and define
\[
 \begin{aligned}
 c&=2+i(T+\tfrac12),& M&=4(A+10)^4,\\
 k&=\min\{j:2^j\ge16M\},& Q&=4k,\\
 \delta&=\frac1{3(Q+2)},&
 L&=\frac1{2^{83}M^{20}[18(Q+2)]^Q}.
 \end{aligned}                                             \tag{RH1}
\]
Every operation in (RH1), including the bounded search for $k$, is primitive
recursive.  Euler--Maclaurin with two correction terms gives
\[
 |\zeta(c+z)|\le M\quad(|z|\le4),\qquad
 |\zeta(c)|\ge\Re\zeta(c)\ge\tfrac14.                    \tag{RH2}
\]
Here the second inequality follows directly from
$\sum_{n\ge2}n^{-2}\le1/4+\int_2^\infty x^{-2}dx=3/4$.
For the first, $\Re(c+z)\ge-2$, $|c+z+j|\le A+10$ for $0\le j\le3$,
and the fourth Bernoulli remainder is bounded by
$|(c+z)_4|/[24(\Re(c+z)+3)]$ after using $|\widetilde B_4|\le1$.
The pole is outside the disk and its distance is bounded below.
These estimates are much smaller than the displayed value of $M$.

Construct a rational complex polynomial $P$ such that
\[
 \sup_{|z|\le4}|P(z)-\zeta(c+z)|\le L/4.                 \tag{RH3}
\]
This is a bounded construction.  One possible recipe uses the Taylor series
at $c$ and the radius-five bound $M_5=8(A+14)^6$, obtained by
Euler--Maclaurin with three correction terms.  Choose $p,b$ with
$2^{-p}\le L/4$ and $M_5\le2^b$, and take
\[
 D=5(p+b+4).
\]
Since $(4/5)^5<1/2$, the Taylor tail beyond degree $D$, on radius four,
is less than $2^{-p}/2$.  Approximating each coefficient to
$2^{-p}/[2(D+1)4^D]$ contributes at most the other half.
Cauchy coefficient integrals can themselves be computed by bounded
quadrature on radius five; the radius-six bound
$M_6=32(A+16)^6$ supplies derivative bounds.  For the $j$th coefficient the
integrand, parametrized on $[0,1]$, has derivative at most
$8(5M_6+jM_5)$ before division by $5^j$.  These explicit bounds give a
primitive recursive sample count.  Thus (RH3) does not invoke an unbounded
search for a successful approximant.

Apply (MM) to $P$.  The possible bad disks have radius $\delta$, and there
are at most $Q$ of them.  Consider the $Q+1$ rational heights
\[
 U_j=T+\frac{j+1}{Q+2},\qquad 0\le j\le Q.              \tag{RH4}
\]
Their spacing is $1/(Q+2)>2\delta$.  Each bad disk can meet at most one of
the horizontal segments $[1/2,2]+iU_j$.  The finite pigeonhole principle
therefore supplies a segment meeting no bad disk.  That entire segment lies
within distance $5/2$ of $c$, so (RH3) and (MM) imply
\[
 |\zeta(\sigma+iU_j)|\ge3L/4\quad(1/2\le\sigma\le2).
                                                               \tag{RH5}
\]

The algorithm need not compute the algebraic roots of $P$.  It simply tests
all the heights (RH4), using a fixed finite mesh.  Let
$J=\lceil12M/L\rceil$ and sample
$\sigma_\ell=1/2+3\ell/(2J)$, $0\le\ell\le J$.
At every sample compute a rational complex value $w_\ell$ with error at most
$L/16$.  Accept a height if
\[
 |w_\ell|^2\ge(5L/8)^2\quad\hbox{for every }\ell.        \tag{RH6}
\]
This is a decidable rational test.  A height satisfying (RH5) passes (RH6).
Conversely, Cauchy's estimate gives $|\zeta'|\le M$ on every candidate
segment, and the mesh spacing is at most $L/(8M)$.  An accepted height has
\[
 |\zeta(\sigma+iU)|\ge7L/16>L/4
 \quad(1/2\le\sigma\le2).                               \tag{RH7}
\]
Return the first accepted height, its mesh, and the margin $L/4$.
There are exactly $Q+1$ candidates and $J+1$ samples per candidate; all
precision and summation bounds have been specified in advance.
This proves a uniform primitive recursive regular-height construction.

\subsection{An explicit certificate for the complete rectangle}
\label{sec:complete-contour}

The accepted horizontal segment is enough to construct every part of the
zero-counting contour with a predetermined precision.  Let its height be
$U\ge8$, and let $L\le1$ be the rational number from the regular-height
construction, so that
$|\zeta(\sigma+iU)|\ge L/4$ for $1/2\le\sigma\le2$.
Put
\[
 \begin{aligned}
 D&=\lceil U\rceil+2,
 & M_\Gamma&=D\,4^{D+D^2},\\
 R&=\lceil U\rceil+3,
 & B&=2(2R+1)R!+R(R+1)(R+1)!,\\
 \mu&=\frac{L}{32M_\Gamma}.
 \end{aligned}                                             \tag{FC1}
\]
The reciprocal-gamma estimate from the analytic foundation gives
$|\Gamma(z)|\ge M_\Gamma^{-1}$ whenever
$\Re z>0$ and $|z|\le D$.
The derivative estimate for the entire completed function gives
$|\xi'(s)|\le B$ on $|s|\le R$.

On the top half-edge $s=\sigma+iU$, $1/2\le\sigma\le2$, we have
$|s(s-1)|\ge1$, $\pi^{-\sigma/2}\ge1/4$, and $|s/2|\le D$.
The defining product therefore gives
\[
 |\xi(s)|\ge \tfrac12\cdot1\cdot\tfrac14
                 \cdot M_\Gamma^{-1}\cdot\tfrac L4=\mu.
\]
On the right edge $s=2+it$, $0\le t\le U$, one has
$|s(s-1)|\ge2$ and $|\zeta(s)|\ge\Re\zeta(s)\ge1/4$.
Consequently $|\xi(s)|\ge1/(16M_\Gamma)\ge\mu$ there.
On the bottom edge $\xi(s)\ge3/8\ge\mu$.
Finally the identity
$\xi(1-\overline s)=\overline{\xi(s)}$
reflects the right edge to the left and the top half-edge to the remaining
top half.  It follows that the complete rectangle has the uniform rational
margin
\[
 |\xi(s)|\ge\mu
 \quad\bigl(s\in\partial([-1,2]+i[0,U])\bigr).          \tag{FC2}
\]
The gamma estimate was used only where $\Re(s/2)>0$; no division by a gamma
value at a pole is hidden in the reflected half.

Divide each horizontal edge into
\[
 J_x=\left\lceil\frac{24B}{\mu}\right\rceil
\]
subintervals and each vertical edge into
\[
 J_y=\left\lceil\frac{8UB}{\mu}\right\rceil
\]
subintervals.  Each segment length is at most $h=\mu/(8B)$.
At every contour vertex $s_j$, evaluate $\xi(s_j)$ by the entire theta
representation to a rational complex number $v_j$ with error at most
$\mu/16$.  The required binary precision is the least $p$ with
$2^{-p}\le\mu/16$; the analytic approximation routine supplies its already
proved primitive recursive tail and quadrature bounds.  Use the same
approximation at the repeated initial/final vertex.

For a contour segment parametrized by
$s(t)=(1-t)s_j+ts_{j+1}$, let
$v(t)=(1-t)v_j+tv_{j+1}$.
The derivative bound implies
\[
 \begin{aligned}
 |\xi(s(t))-v(t)|
 &\le (1-t)Bt|s_{j+1}-s_j|
       +tB(1-t)|s_{j+1}-s_j|+\mu/16\\
 &\le Bh/2+\mu/16\le\mu/8.
 \end{aligned}
\]
The straight homotopy between the analytic image and this rational polygon
therefore has modulus at least $7\mu/8$ everywhere.
The polygon has exactly $2(J_x+J_y)$ segments; its winding integer is
computed by finite rational arithmetic and equals the analytic contour
count.

For example, no trigonometric or branch computation is needed to find that
integer.  For each oriented edge from $a$ to $b$, add one if
$\Im a\le0<\Im b$ and $\det(a,b)>0$; subtract one if
$\Im b\le0<\Im a$ and $\det(a,b)<0$; add zero otherwise.
Here $\det(a,b)=\Re a\Im b-\Im a\Re b$.
The half-open conventions handle vertices on the real axis.
If an edge straddles that axis and its determinant vanishes, it would pass
through the origin, which (FC2) and the homotopy bound exclude.
Thus all branches are decidable exact rational comparisons.
This completes the primitive recursive contour constructor, rather than
merely assuming the existence of the other three zero-free edges.

\section{The zero-counting argument}
\label{sec:counting}

Let
\[
 \xi(s)=\tfrac12s(s-1)\pi^{-s/2}\Gamma(s/2)\zeta(s).
\]
The analytic foundation supplies its entire continuation and the identities
$\xi(1-s)=\xi(s)$ and $\xi(\overline s)=\overline{\xi(s)}$.
For a rational regular height $U\ge8$ equipped with a positive margin on
$[1/2,2]+iU$, let $N(U)$ denote the winding count of $\xi$ on the
positively oriented rectangle $[-1,2]+i[0,U]$.  Its classical interpretation
is the number of nontrivial zeros with $0<\Im\rho<U$, counted with
multiplicity.  The supplied margin, the functional equation, and the
explicit nonvanishing bounds on the other edges turn this contour into a
finite rational winding certificate.

\subsection{The exact argument identity}

The function $\xi$ is positive on the real interval $[-1,2]$.
An explicit lower bound $1/4$ follows from the theta integral
\[
 \xi(s)=\frac12+\frac{s(s-1)}2
 \int_1^\infty \Bigl(\sum_{n\ge1}e^{-\pi n^2x}\Bigr)
       (x^{s/2}+x^{(1-s)/2})\,\frac{dx}{x}.
\]
For $s\in[0,1]$ the integral is at most one: use
$\sum_{n\ge1}e^{-\pi n^2x}\le2e^{-3x}$ and
$x^{s/2}+x^{(1-s)/2}\le2x$.
Thus $\xi(s)\ge3/8$ there; outside $[0,1]$ the integral term is
nonnegative.

Start the continuous argument at zero at $s=2$, and follow the path
$2\to2+iU\to1/2+iU$.  Reflection
$s\mapsto1-\overline s$ and conjugation show that the change of argument
on the whole rectangle is twice the change on this path.  The latter is
\[
 \pi+\vartheta(U)+\arg\zeta(1/2+iU),
\]
where
\[
 \vartheta(U)=\Im\log\Gamma(1/4+iU/2)-\tfrac U2\log\pi
\]
uses the logarithm of gamma continuous in the right half-plane, real on the
positive real axis.  In particular, it is not a repeatedly reset principal
argument.  The factor $s(s-1)$ contributes exactly $\pi$: at the endpoint
its value is $-U^2-1/4$, and both individual arguments were continued through
the upper half-plane.  Therefore
\[
 N(U)=1+\frac{\vartheta(U)}\pi+S(U),\qquad
 S(U)=\frac1\pi\arg\zeta(1/2+iU).                       \tag{C1}
\]
All argument identities can equivalently be read as identities of the
integer winding numbers of rational polygonal approximations, whose
homotopies have the supplied positive margins.

\subsection{Backlund's estimate with a finite polynomial crossing count}

Euler summation gives, for $\Re s>0$,
\[
 \zeta(s)=\frac{s}{s-1}-s\int_1^\infty\{x\}x^{-s-1}\,dx.
\]
Consequently, if $U\ge8$ and $|z-2|\le7/4$, then
\[
 |\zeta(z\pm iU)|\le5(U+4).
\]
Indeed $\Re z\ge1/4$, $|z\pm iU|\le U+4$, and
$|z\pm iU-1|\ge U-7/4>2$.
At $z=2$, the real part of $\zeta(2+iU)$ is at least $1/4$.

Approximate $w\mapsto\zeta(2+w+iU)$ by a rational polynomial $P(w)$
on $|w|\le7/4$, with error less than both $1/8$ and one quarter of the
supplied horizontal nonvanishing margin.  Let
\[
 Q(w)=\tfrac12\bigl(P(w)+\overline{P(\overline w)}\bigr).
\]
This is a real polynomial, $Q(0)\ge1/8$, and
$|Q(w)|\le6(U+4)$ on the outer disk.  By (J), the number $q$ of its
roots in $|w|\le3/2$, counted with multiplicity, satisfies
\[
 q\log(7/6)\le\log(48(U+4)),\qquad
 q\le7\log(48(U+4)),                                   \tag{C2}
\]
because $\log(7/6)\ge1/7$.
The real roots of $Q$ partition the segment $[-3/2,0]$ into at most
$q+1$ intervals.  On each interval $P$ stays in one open half-plane, so its
continuous argument changes by less than $\pi$.
Homotopy from $P$ to $\zeta$ changes only the endpoint arguments, by a
total of less than $\pi$; the vertical path at real part two stays in the
right half-plane.  Allowing an extra unit for all endpoint conventions gives
\[
 |S(U)|\le7\log(48(U+4))+4.                             \tag{C3}
\]
This proof uses only the finitely many roots of a rational real polynomial,
not a procedure deciding the zeros of the real part of an analytic
function.  Multiple roots cause no problem and only enlarge the estimate.

\subsection{Riemann--von Mangoldt with an explicit error}

The explicit complex Stirling estimate proved in the analytic foundation is
\[
 \left|\vartheta(U)-\left(\frac U2\log\frac U{2\pi}
                    -\frac U2-\frac\pi8\right)\right|
 \le1\qquad(U\ge8).
\]
Combining this with (C1)--(C3), and using $\log48<4$, gives the deliberately
coarse but uniform statement
\[
 \boxed{\quad
 \left|N(U)-\frac U{2\pi}\log\frac U{2\pi}
                 +\frac U{2\pi}\right|
 \le8\log(U+4)+40.
 \quad}                                                 \tag{RVM}
\]
The quantifiers are important: $U$ is supplied with a finite nonvanishing
certificate, or is returned by the bounded regular-height algorithm.
No exact decision whether an arbitrary real height is a zero ordinate is
used.

For an integer $n\ge8$, let $U_n\in(n,n+1)$ be the deterministically
selected height and let $C(n)=N(U_n)$.  The contour construction and winding
computation make $C$ a total primitive recursive integer function.
Moreover $C(n+1)\ge C(n)$ and, with
$\mathcal M(t)=t\log(t/(2\pi))/(2\pi)-t/(2\pi)$,
\[
 |C(n)-\mathcal M(n)|\le9\log(n+5)+40.                           \tag{RVM-PR}
\]
This is one particularly simple arithmetical formulation of the classical
asymptotic.  In the classical interpretation, zero counts at any intermediate
height are bracketed by the adjacent certified counts.  A finite-precision
constructive statement uses those brackets rather than a discontinuous
real-to-integer counting function.


\section{Distribution consequences}
\label{sec:consequences}

\subsection{Local counts and multiplicities}

All zeros in the slab $0<\Re s<1$, $T\le\Im s\le T+1$ lie in the disk
of radius $5/2$ around $2+i(T+1/2)$.  Apply (J) with outer radius four and
inner radius $5/2$.  Approximating $\zeta$ by a rational polynomial closely
enough to preserve any supplied finite family of zero-cluster certificates
gives
\[
 \sum_j \operatorname{count}(\mathcal B_j)
 \le\frac{\log(16M)}{\log(8/5)}
 \le12\log(T+11)+12,\qquad T\ge8.                       \tag{LC}
\]
Here the $\mathcal B_j$ are disjoint certified clusters in the slab and
$M=4(\lceil T\rceil+10)^4$.
The inequalities $\log(8/5)\ge3/8$ and $\log64<9/2$ suffice for the last
estimate.  This proves the usual $O(\log T)$ local zero-count bound and the
same bound on the multiplicity of an individual zero.  The certificate
form does not require a decision procedure for exact equality of ordinates.

\subsection{Infinitely many zeros, with a bounded witness}

Take an integer $T\ge2^{16}$ and choose regular heights
$U\in(T,T+1)$ and $V\in(2T,2T+1)$.  For $x\in[T+1,2T]$,
\[
 \mathcal M'(x)=\frac1{2\pi}\log\frac{x}{2\pi}
       \ge\frac1{16}\log T,
\]
since $2\pi<8$, $\log(2\pi)<3$, and $\log T\ge8$.
Thus $\mathcal M(V)-\mathcal M(U)\ge T\log T/32$.
The sum of the two errors in (RVM) is at most $30\log T$.
As $T\ge2^{16}>1920$,
\[
 N(V)-N(U)\ge\frac{T\log T}{64}\ge\frac T8.            \tag{AB}
\]
The two contour certificates therefore give a positive-count rectangle
strictly above height $T$, and indeed a count at least $T/8$.
Using $T=4^{m+8}$ gives disjoint such height ranges for successive $m$.
This is a bounded constructive witness for unboundedly many nontrivial
zero clusters in disjoint regions. The following construction extracts
a primitive recursive point name from each such cluster, yielding
a uniform sequence of pairwise apart nontrivial zeros. It does not
order all exact zero ordinates.

\subsection{Extracting an actual zero name from a positive count}
\label{sec:zero-extraction}

A positive contour count contains more constructive information than an
unrelated sequence of small residuals. Here is a coherent root extraction,
using only the quantitative minimum-modulus construction and the finite
winding rules already proved.

Fix one rational rectangle $Q_0$ with a certified zero-free boundary and
positive xi winding. Choose an integer $A\ge1$ so that
$|s-2|\le A$ throughout $Q_0$, and use $f(z)=\xi(2+Az)$. The entire
theta-integral construction gives, with $R=4A+2$,
\[
        |f(z)|\le M:=1+2R(R+1)R!\quad(|z|\le4),
             \qquad f(0)=\xi(2)=\zeta(2)/\pi\ge1/4.
\]
The same expression with $R=5A+2$ bounds the larger disk used for Taylor
approximation. The minimum-modulus lemma therefore applies with the fixed
integer $K=4\lceil\log_2(16M)\rceil$. In its notation one may take
\[
       L_\delta=2^{-83}M^{-20}(\delta/6)^K,
\]
so that outside at most $K$ disks of radius $\delta$ the normalized
function has modulus at least $3L_\delta/4$. The global derivative on
$Q_0$ is bounded by a fixed primitive recursive integer $D\ge1$; for
example the Cauchy bound from the displayed disk estimate permits $D=M$.
These constants may be very large, but none is obtained by an unbounded
search.

For a current rectangle $[a,a+w]\times[b,b+h]$, use the $K+1$ candidates
\[
 x_j=a+w/3+\frac{(j+1)w}{3(K+2)},\qquad
 y_j=b+h/3+\frac{(j+1)h}{3(K+2)}\quad(0\le j\le K),
\]
and put $\delta=\min(w,h)/(9A(K+2))$. Twice this radius is smaller
than either candidate spacing in normalized coordinates, so each exceptional
disk excludes at most one vertical candidate and at most one horizontal
candidate. There is therefore a safe full segment in each direction.
On every candidate segment evaluate $\xi$ to error $L_\delta/16$ on a
mesh of physical width at most $L_\delta/(8D)$. Choose the first segment
whose rational samples have modulus at least $L_\delta/2$, tested by
squaring. A safe candidate passes this test, and every passing candidate
has true modulus at least $5L_\delta/16$ on the entire segment. Record
the smaller margin $L_\delta/4$ for each of the two chosen cuts.

If the parent's recorded margin is $m$, give every child the common margin
\[
                         m'=\tfrac12\min(m,L_\delta/4).
\]
Every child edge then has true
modulus at least $2m'$. Restrict the inherited boundary certificates
to the child edges, and refine both inherited edges and the two cuts to
a common mesh of width at most $m'/(8D)$, with sample error at most $m'/8$.
The resulting Lipschitz certificates verify the claimed margin $m'$.
Reuse
identical samples with opposite orientations on shared child edges. This
constructs full closed-boundary certificates for all four children, not
just certificates for their new cuts. Each side length lies between
$1/3$ and $2/3$ of the corresponding parent side length.

Each child winding is nonnegative. To see this without already assuming
analytic root existence, approximate $\xi$ on that child's boundary by a
rational complex polynomial to less than half the certified boundary
margin. The finite zero-free homotopy identifies the two windings. The
polynomial engine of Section~\ref{sec:polynomial-algebra} factors that
polynomial, so its winding is the sum of the nonnegative indices of its
linear factors. Cancellation of the new internal edges makes the four
child windings sum to the parent's winding. At least one is consequently
positive; choose the first such child in a fixed ordering.

Iterating this procedure gives nested rational rectangles $Q_j$ with
positive winding and diameter at most $D_0(2/3)^j$, for a supplied
$D_0\ge1$. All searches are bounded, and the bounds are recomputed by
primitive recursion from the current rational rectangle and its inherited
boundary margin. In particular this is not a search that waits until a
possibly vanishing sample becomes visibly nonzero.

Let $c_j$ be the centre of $Q_j$. Then
\[
                        |\xi(c_j)|\le D\,\operatorname{diam}(Q_j).
\]
Indeed, a strict violation would put the entire boundary image in a disk
centred at $\xi(c_j)$ which misses zero; its winding would be zero. The
finite-precision version proves the displayed weak inequality by ruling
out every rationally witnessed violation, using the same derivative bound.

Choose an integer $L$ with $2^L\ge4DD_0$, and set
\[
                         \alpha(k)=c_{\,2(k+L)}.
\]
Since $(2/3)^2<1/2$, nestedness gives, for $\ell\ge k$,
\[
       |\alpha(\ell)-\alpha(k)|\le2^{-k-2},\qquad
       |\xi(\alpha(k))|\le2^{-k-2}.
\]
The centres therefore form one primitive recursive fast Cauchy name for an
actual zero. The initial positive boundary margin and derivative bound
also keep this zero positively separated from $\partial Q_0$, so it lies
inside the original rectangle. Every recurrence invariant used to obtain these inequalities
concerns finite rational rectangles, contour certificates, signed integer
counts and fixed-precision approximation errors. No countable choice or
decision of exact zero ordinates is involved.

For rectangles above the real axis, the explicitly nonvanishing factors
relating $\xi$ and $\zeta$ identify the extracted point as a nontrivial
zeta zero. Applying the construction to disjoint certified height bands
gives pairwise apart zero names. This proves constructive infinitude when
the counting estimate makes those band counts positive; it does not supply
a procedure for deciding equal ordinates or sorting all exact zeros.


\subsection{Average density in growing windows}

From (RVM), for $T\le U<T+1$ and
$T+H\le V<T+H+1$, where $1\le H\le T$, one obtains
\[
 N(V)-N(U)=\frac{H}{2\pi}\log\frac T{2\pi}
          +O\left(\frac{H^2}{T}+\log(T+H+4)\right).
\]
The constants here can be read directly from (RVM) and the derivative of
$\mathcal M$; alternatively one retains the exact difference $\mathcal M(V)-\mathcal M(U)$ and the
explicit error sum.  Hence when $H\to\infty$ and $H=o(T)$, the mean density
in such a window is asymptotic to $\log T/(2\pi)$.
For a unit window the error is the same order as the main term; the correct
consequence of this proof is only $N(T+1)-N(T)=O(\log T)$, not an
asymptotic equivalent.  Similarly, the usual inverse relation
$\gamma_n\sim2\pi n/\log n$ is a classical consequence of the count
asymptotic.  Its constructive reading is via inverse count brackets and
does not silently posit a uniformly decidable ordering of equal zero
ordinates.

\subsection{An explicit growing-window estimate}
The asymptotic just stated can be given a convenient numerical form.
Let $T\ge8$, $1\le H\le T$, and choose the certified heights
$U\in(T,T+1)$ and $V\in(T+H,T+H+1)$. By
\eqref{eq:main-rvm},
\begin{equation}
 \left|N(V)-N(U)-\frac{H}{2\pi}\log\frac{T}{2\pi}\right|
 \le\frac{H^2}{12T}+18\log(T+H+5)+80.
 \label{eq:window-explicit}
\end{equation}
Indeed the change from $U,V$ to $T,T+H$ in the main-term difference
is at most $2\log(T+H+5)$. The change of the logarithm inside
$\int_T^{T+H}\mathcal M'(x)dx$ contributes at most
$H^2/(4\pi T)\le H^2/(12T)$. Finally the two counting errors are
at most $16\log(T+H+5)+80$.

For a particularly simple arithmetic specialization take $T=n^2$,
$H=n$, with integer $n\ge3$. The right side becomes
$1/12+18\log(n^2+n+5)+80$, while the leading term is
$n\log(n^2/(2\pi))/(2\pi)$. Dividing proves the claimed mean density
with a primitive recursive convergence modulus. One need not introduce
an irrational square-root input into the checker.

\subsection{Effective tails of reciprocal-zero sums}
\label{sec:zero-sums}

For every certified height $u\ge8$, the coarse consequence
\begin{equation}
                    N(u)\le64u\log(u+2)
 \label{eq:coarse-count}
\end{equation}
follows immediately from~\eqref{eq:main-rvm}. For example, the main
term is at most $u\log(u+2)/6$, and its remaining error is at most
$36\log(u+2)\le(9u/2)\log(u+2)$. The factor $64$ leaves ample
slack for rational endpoint adjustments.

\begin{theorem}[Uniform reciprocal-square tail]
For every integer $T\ge8$, every finite collection of zeros with
$|\Im\rho|>T$, counted with multiplicity, satisfies
\begin{equation}
       \sum |\rho|^{-2}
       \le\frac{2048(\log(T+2)+1)}{T}.
 \label{eq:reciprocal-tail}
\end{equation}
The finite constructive formulation uses disjoint certified clusters
and the lower bound $|\rho|\ge|\Im\rho|$ in each height band.
\end{theorem}

\begin{proof}
Partition positive ordinates into the dyadic bands
$(2^jT,2^{j+1}T]$. To include a possible zero at the upper endpoint,
choose a certified height $u_j\in(2^{j+1}T,2^{j+1}T+1)$.
By~\eqref{eq:coarse-count}, the count of positive ordinates in that
band is at most
\[
 144\,2^jT\bigl(\log(T+2)+(j+1)\log2\bigr).
\]
Here $2^{j+1}T+1\le(9/8)2^{j+1}T$ and
$2^{j+1}T+3\le2^{j+1}(T+2)$ suffice. Conjugation doubles this
bound for both signs. Division by $(2^jT)^2$ and summation give
\[
 \frac{288}{T}\sum_{j\ge0}2^{-j}
       \bigl(\log(T+2)+(j+1)\log2\bigr)
 =\frac{576\log(T+2)+1152\log2}{T},
\]
which is smaller than~\eqref{eq:reciprocal-tail}.
The actual arithmetic argument uses a finite geometric sum and its
explicit remainder. It proves the same bound for every finite number
of bands, so no induction over a quantified convergence assertion is
being used.
\end{proof}

It follows that $\sum_\rho\rho^{-m}$ is absolutely convergent for
each integer $m\ge2$. For $m=1$, pair each positive-height zero
$\rho$ with the negative-height zero $1-\rho$. The identity
\begin{equation}
 \frac1\rho+\frac1{1-\rho}
       =\frac1{\rho(1-\rho)}
 \label{eq:paired-reciprocals}
\end{equation}
gives modulus at most $1/(\Im\rho)^2$. Hence this paired reciprocal
series is absolutely convergent as a series of pairs, without a
Riemann-hypothesis assumption. This is not a claim of absolute
convergence of the unpaired series $\sum_\rho1/\rho$.

An explicit uniform tail choice for either assertion is
\begin{equation}
                         T_p=2^{2p+32}.
 \label{eq:reciprocal-modulus}
\end{equation}
Since $\log(T_p+2)+1\le2p+34$ and
$p+17\le2^{p+5}$, the right side of
\eqref{eq:reciprocal-tail} is less than $2^{-p}$.

The finite partial sums also have a constructive representation which
does not list exact zeros. Put $\eta=1/4096$. The derivative bound for
$\xi$ with $R=3$ is $D_3=372$, so its positive real-axis bound gives
$|\xi(x+\ii\eta)|>1/4$ for $-1\le x\le2$.
On the certified rectangle $[-1,2]+\ii[\eta,U]$, integrate
\[
 \frac{1}{2\pi\ii}\frac{\xi'(s)}{\xi(s)}w(s),
 \quad w(s)=s^{-m}\quad\hbox{or}\quad w(s)=\frac1{s(1-s)}.
\]
The weights have no pole in this rectangle. Division has a positive
boundary clearance; derivatives and quadrature have the explicit bounds
already established. Polynomial factorization and the elementary
one-factor residue calculation identify this integral with the
corresponding finite zero sum, counting multiplicity. For $m\ge2$
include the reflected lower rectangle as well; for the paired sum the
upper rectangle alone suffices. Approximate these contour integrals at
precision $p+2$ and use~\eqref{eq:reciprocal-modulus} for the tail.
This supplies primitive recursive Cauchy names for the stated sums.

\section{An exact-rational certificate for a simple critical-line zero}
\label{sec:example}

The preceding general contour bounds establish totality but are far too
large for a practical run. The archive therefore also contains a small
independent calculation using much sharper local information. This is a
worked arithmetic certificate, with its analytic interpretation proved
below. It illustrates the distinction between executing a calculation
and formalizing that calculation's soundness in the object theory.

Take the exact rational data
\begin{equation}
 c=\frac12+14.1347251417347\,\ii,
 \qquad r=h=10^{-6},\qquad B_2=4096.
 \label{eq:example-data}
\end{equation}
The decimal in this display is a terminating rational, not an assumed
exact zeta ordinate. Use~\eqref{eq:em} with $N=64$ and $m=10$ to
enclose $f(c)$, $f(c-h)$ and $f(c+h)$, where $f=\zeta$. The
implementation uses exact rational arithmetic, 384-bit dyadic centres,
outward-rounded radii, the atanh logarithm series, and the exponential
Taylor series. The Bernoulli numbers are computed by their rational
recurrence. The remainder is bounded by
\eqref{eq:em-bound}, multiplied by two for the $\normone{\cdot}$ radius.

Let $d$ be the rational centre of
$(f(c+h)-f(c-h))/(2h)$, and let $e_d$ be its evaluation error plus
$B_2h/2$. Then
\begin{equation}
 |f'(c)-d|\le e_d.
 \label{eq:example-derivative}
\end{equation}
To see this, integrate $f'$ over the real segment $[c-h,c+h]$.
The mean error from replacing $f'(c+x)$ by $f'(c)$ is at most
$(2h)^{-1}\int_{-h}^h B_2|x|\,dx=B_2h/2$.

Here is a completely elementary justification of the second-derivative
bound. On the disk and the finite-difference segment,
$\sigma\ge1/4$, $|s|\le16$ and $|s-1|\ge1$.
Differentiating~\eqref{eq:euler-basic} twice and using
\eqref{eq:logmoments} gives
\[
 |\zeta''(s)|\le 2+2\cdot16+16\cdot128=2082<4096.
\]
The same inequalities hold throughout a neighborhood of these sets,
as can be checked directly from~\eqref{eq:example-data}.

The script verifies, as exact rational comparisons,
\begin{align}
 A_0+e_dr+\tfrac12B_2r^2&< |d|_\infty r,
       \label{eq:example-rouche}\\
 |d|_\infty-e_d-B_2r&>0,
       \label{eq:example-simple}
\end{align}
where $A_0$ is the certified upper bound for $|f(c)|$.
Some outward numerical summaries of the exact inequalities are
\begin{center}
\begin{tabular}{ll}
\toprule
Quantity & Conservative decimal summary\\
\midrule
$A_0$ & less than $5.639\times10^{-15}$\\
$\Re d$ & between $0.78329651186710$ and $0.78329651186712$\\
$\Im d$ & between $0.12469982974821$ and $0.12469982974823$\\
Left side of~\eqref{eq:example-rouche}
  & less than $4.097\times10^{-9}$\\
Right side of~\eqref{eq:example-rouche}
  & greater than $7.832\times10^{-7}$\\
Lower bound in~\eqref{eq:example-simple}
  & greater than $0.77715$\\
\bottomrule
\end{tabular}
\end{center}
The JSON file retains the exact rational values; these displayed
decimals are not used to make the decisions.

Taylor's formula with the explicit integral remainder gives, on
$|z-c|=r$,
\[
 |f(z)-d(z-c)|\le A_0+e_dr+B_2r^2/2<|d(z-c)|.
\]
Rouch\'e stability therefore gives count one in the disk. Also
$|f'(z)|\ge |d|_\infty-e_d-B_2r>0$ throughout the disk, so its
zero is simple.

\subsection{A primitive recursive point name}
There is an even more explicit constructive existence argument here.
Use the fixed rational $d\ne0$ and the map
\[
             \Phi(z)=z-f(z)/d.
\]
The rational certificate checks
\begin{equation}
 \kappa:=\frac{e_d+B_2r}{|d|_\infty}<\frac1{64},\qquad
 \frac{A_0}{|d|_\infty}<\frac r{64}.
 \label{eq:example-contraction}
\end{equation}
Thus $\Phi$ maps the closed disk into the concentric disk of radius
$r/32$ and is a contraction of factor at most $1/64$.
Starting at $c$, $n$ exact iterations have successive errors bounded
geometrically, so their distance from the root is at most
$2r64^{-n}$. For a requested precision, perform a predetermined
number of iterations and evaluate each operation with a predetermined
rounding budget; the geometric stability bound controls the sum of
rounding errors. For example, use $n=p+30$ iterations and a per-step
error at most $2^{-p-40}r$, increasing the elementary working precision
by the explicit division bound $|d|>1/2$. The final error is less
than $2^{-p}$. This is a primitive recursive Cauchy name. The same
contraction estimate proves uniqueness.

The disk is invariant under $z\mapsto1-\bar z$. The completed
functional equation sends its root to another root in the same disk;
uniqueness implies $\rho=1-\bar\rho$, hence $\Re\rho=1/2$.
The certificate therefore proves one simple critical-line zero in
the disk~\eqref{eq:example-data}. It does not claim that this is the
lowest positive zero; that extra statement would require a count in
the entire lower region.

\subsection{What was executed}
\code{python3 code/rational_zeta.py} recomputes all three zeta
enclosures and all displayed strict rational tests. It writes
\code{verification/simple_zero.json}. No floating-point special-function
evaluation, zero table or root finder is used by this program.
Its correctness is justified by the mathematical estimates just given.
The program is not a proof checker for full $\ipra$, and its successful
execution is not reported as a Lean kernel proof.

\section{Concrete finite arithmetic sentences}
\label{sec:contracts}

This section fixes representative internal statements rather than leaving
``formalize the analytic theorem'' as an unspecified goal. Integer and
rational coding conventions can be chosen freely, provided their
operations are compiled to the repository's primitive-recursive syntax.
Invalid numerical codes are assigned fixed default values. A conditional
Boolean checker returns one on an invalid hypothesis code; a constructor's
separate validity theorem ensures that its own outputs are genuine
certificates.

\subsection{The primary counting sentence}
Fix the deterministic height selector and contour constructor of
Section~\ref{sec:regular-heights}. For integer $n\ge8$ let
$\mathsf{Gen}(n)$ be its complete rational certificate and
$C(n)$ its crossing-sum integer. Return zero for $n<8$ in the
definition of $C$. Let $\mathsf{Valid}$ recompute the specified xi
approximants, verify the contour mesh, prescribed error bounds and
positive sample margins, and recompute the crossing sum. Two separate
internal theorems are required:
\begin{equation}
 \ipra\vdash\forall n\ge8\quad
            \mathsf{Valid}(\mathsf{Gen}(n))=1,
 \label{eq:generator-sentence}
\end{equation}
and the numerical estimate below. The first is supplied by the
minimum-modulus and finite pigeonhole argument, including the other
three edges of the rectangle.

Let $m(n,p)$ and $\ell(n,p)$ be the fixed rational approximants, with
error at most $2^{-p}$, to
\[
 \mathcal M(n)=\frac{n}{2\pi}\log\frac{n}{2\pi}-\frac{n}{2\pi},
 \qquad \log(n+5),
\]
respectively. Their definitions use the log and exponential routines
and extra working precision sufficient for multiplication by $n$ and
division by $2\pi\ge6$. Define a Boolean function
$\mathsf{CountCheck}(n,p)$ to return one if $n<8$, and otherwise to
test the exact rational inequality
\begin{equation}
 |C(n)-m(n,p)|\le9\ell(n,p)+40+10\,2^{-p}.
 \label{eq:count-matrix}
\end{equation}
Then the intended object sentence is precisely
\begin{equation}
       \ipra\vdash\forall n\,\forall p\quad
                   \mathsf{CountCheck}(n,p)=1.
 \label{eq:count-sentence}
\end{equation}
The $10\,2^{-p}$ budget is the sum of the error in $m$ and nine
times the error in $\ell$. The proof of (RVM-PR) establishes this
inequality with $n,p$ free. For a supplied regular contour, the analogous
checker uses $8\log(U+4)+40$ and a $9\,2^{-p}$ budget.

\subsection{Representative zero-exclusion and identity sentences}
For integer $H\ge1$ and a rational input $s$ in the rectangle
$|\Im s|\le H$, $|\Re s-1|\le\varepsilon_H$, the lower bound
$b_H$ in Theorem~\ref{thm:elem-F-clearance} is rational. Compute the
symbolically pole-removed numerator $F$ with error at most $b_H/8$.
The returned rational $v$ must satisfy
\[
                         |v|^2\ge(7b_H/8)^2.
\]
The checker tests the input inequalities, recomputes $v$, and performs
this squared-norm comparison. The working precision is the bounded
binary logarithm of $8/b_H$ plus the stated arithmetic error budget.
Thus this is one uniform finite test at an explicit precision, rather
than a negation of an undecidable equality.

For $\xi$ symmetry, use one common box and the same theta summation
parameters at $s,1-s,\bar s$. The rational expression for the completion
then satisfies the reflection identity exactly. The conjugation identity
is also exact if the elementary evaluator treats conjugate inputs by
conjugating its rational Taylor polynomial and the chosen real constants.
The Mellin identification remains a separate, substantive finite theorem.

For trivial zeros and their derivative formula, compile the finite gamma
shift identity and the Taylor coefficient identity at the integer input
$-2n$. At precision $p$ the equality checker compares the two prescribed
rational approximants with the sum of their proved error budgets. The
strict simplicity assertion uses the rational lower bound
$(2n)!/(2\,8^{2n})$ to select a fixed successful derivative precision.
No equality test on an arbitrary real or complex input is part of this
checker.

\subsection{Reciprocal sums as Cauchy names}
For each integer $m\ge2$, define $A_m(p)$ by a weighted finite contour
integral at a selected regular height above $T_{p+3}$, using both
upper and lower rectangles. Compute the quadrature and function values
to total Euclidean error at most $2^{-p-3}$. The tail has the same
bound by~\eqref{eq:reciprocal-modulus}. After conversion to the
rational L1 norm, the error is at most $2^{-p-1}$. The internal
statement can therefore be the exact rational inequality
\begin{equation}
 \normone{A_m(p)-A_m(q)}\le2^{-p}+2^{-q}
       \qquad(m\ge2).
 \label{eq:cauchy-sentence}
\end{equation}
All loop bounds remain primitive recursive when $m$ is a numerical
parameter: the weight and its derivatives involve finite powers, and
the lower contour has the fixed rational height $\eta=1/4096$.
For the paired reciprocal sum use its upper-contour weight
$1/[s(1-s)]$ instead. The same sentence and budgets apply.

\subsection{Why these are uniform proofs}
The proof uses several nested loops, but every nesting has a numerical
bound supplied before execution. Examples are the degree-bounded gcd,
finite determinant permutation list, bounded algebraic root grid,
fixed-precision bisections, Euler--Maclaurin truncations, candidate-height
list, contour mesh and crossing sum. Their state is a finite numerical
array. At each loop, all later precision demands are already numerical
parameters. Identities and rational error inequalities about that state
are quantifier-free invariants.

For a bounded family of such inequalities, use the Boolean product of
their characteristic functions. Induction applies to the equation that
this product equals one, with the bound as a free parameter. It does
not apply to a statement quantifying over every precision or asserting
existence of an unspecified real object. Polynomial root names in the
finite algebraic intermediate calculus have one fixed bounded
evaluation routine on their finite codes, so no universal evaluator for
arbitrary primitive recursive programs is being assumed.

These details specify the arithmetic claims and the human proofs of
their uniform bounds. The remaining Lean implementation work is to
compile the fixed programs, prove the indicated invariants in the deeply
embedded theory, and build the actual derivations of
\eqref{eq:generator-sentence}, \eqref{eq:count-sentence}, and
\eqref{eq:cauchy-sentence}. A theorem about their truth in ambient Lean's
standard natural numbers is a different endpoint.

% Section source for the main manuscript; uses standard AMS math commands.
\section{The exact arithmetic foundation and the Lean endpoint}
\label{sec:foundation-audit}

\subsection{What the repository supplies}
We inspected ProveIt at the following commit:
\begin{center}\small\ttfamily
b5a3ad037723c44d69a3d5b28a950c554327a989.
\end{center}
The file \texttt{ConciseIPRA/PR.lean} defines a deeply embedded first-order
theory: its function symbols are syntax trees built from zero, successor,
projections, composition, and primitive recursion. Its nonlogical axioms
are the defining equations, separation and injectivity of successor, and
induction for quantifier-free formulas. The deduction rules in
\texttt{Deduction.lean} are intuitionistic natural deduction with equality.
The notation
\[
       \Gamma\vdash[\mathsf{PRAAx}]\varphi
\]
denotes a derivation in that object theory; it does not merely assert that
the standard natural numbers satisfy~$\varphi$.

The native scalar calculus uses iterator terms and formulas consisting
only of equations and two binding constructions. Its ordinary readings are
\[
 Q_yB(y)\;\equiv\;B(0)\to\forall y\,B(y),\qquad
 D_{x,y}B(x,y)\;\equiv\;
     \exists x\bigl(B(x,0)\to\forall y\,B(x,y)\bigr).
\]
The inspected theorem \texttt{pra\_iff\_native}, in
\texttt{NativeBiInterpretation.lean}, says exactly that
\[
 (\Gamma\vdash[\mathsf{PRAAx}]\varphi)
 \quad\Longleftrightarrow\quad
 \mathsf{NPrf}\bigl(\Gamma.\mathsf{map}\,\mathsf{NFm.ofPRA}\bigr)
                   (\mathsf{NFm.ofPRA}\,\varphi).
\]
Both interpretations preserve the whole numerical domain and the
repository also proves both formula round trips. Consequently the natural
formalization route is to prove the finite arithmetic estimates in the
ordinary iPRA presentation, then apply this established bridge. Rebuilding
the smallest native proof checker is not a prerequisite for this work.

\subsection{Finite statements rather than hidden higher-order quantifiers}
Every rational is coded by a signed numerator and a positive denominator;
complex rationals are pairs of rationals. Finite lists, polynomial
coefficients, meshes, contour samples, and finite certificates are numerical
codes. Their comparisons are decidable by integer arithmetic. Malformed
codes receive a fixed default output, so each implementation is a total
primitive recursive function. Formula quantifiers range over such natural
codes, not over arbitrary real numbers, analytic functions, or primitive
recursive programs.

For a fixed Boolean checker $C$ with $r$ numerical inputs, its internal
correctness statement is
\[
         \forall x_1\cdots\forall x_r\;C(x_1,\ldots,x_r)=1.
\]
The quantifier-free matrix is an equation in the language of iPRA. A
conditional checker returns $1$ when its finite input hypotheses are false
and otherwise checks the asserted rational conclusion. This convention
does not dispense with hypotheses: the report must specify exactly which
hypotheses the finite checker recognizes.

Statements for a primitive recursive Cauchy name are schemas in a
\emph{fixed definition} of that name. Its regularity hypotheses may remain
universally quantified assumptions. No step evaluates all primitive
recursive program codes by a single total primitive recursive function.
That distinction is essential: a universal evaluator of that kind is not
available.

At each use of induction the requested precision and all size parameters
are fixed first. The invariant is an equation or inequality about rational
partial sums, products, mesh-cell errors, or a bounded Boolean fold. A
bounded conjunction $\bigwedge_{j<n}P(j)$ is represented by the product of
the characteristic values of the fixed predicate~$P$. An induction over
such a fold is quantifier-free induction. This does not authorize induction
over an assertion such as ``the analytic equality holds at every
precision'', or over the existence of an unspecified zero. Those larger
claims are assembled only after the finite inductions have finished.

\subsection{What a contour certificate certifies}
A finite contour certificate records a rational polygon in the domain,
rational samples of the fixed analytic approximation, a numerical sample
precision, a derivative or continuity bound, a positive rational margin,
and an integer winding count. Its checker recomputes the samples and tests
the error and separation inequalities. Sample values are not trusted
merely because the certificate contains them.

For example, if the approximation error is at most $e$, each domain edge
has length at most $h$, and the derivative bound is~$L$, then an inequality
of the form
\[
            |w_j|_\infty>4e+4Lh
\]
at every sample gives a zero-free homotopy from the actual boundary image
to its polygonal approximation, after including the corresponding
successive-sample bounds. All tests can use rational norms; no square-root
oracle is required. The winding of the rational polygon is computed by
oriented ray crossings. A transverse rational ray can be selected by a
finite search which excludes the finitely many vertex directions.

There are two separate formal obligations. First, a valid contour has the
claimed count and obeys the distribution estimate. Second, the specified
primitive recursive construction produces valid contours at the stated
heights. The first obligation alone is insufficient for unconditional
unboundedness of the zero count. The minimum-modulus and regular-height
construction in the counting argument exists precisely to supply the
second obligation. Nor is a certified count the same as an algorithm that
sorts exact zeros by ordinate or decides equal ordinates.

\subsection{The inspected analytic drafts and their status}
The repository's archived functional-equation draft supplies a quantitative
human proof of the theta transformation, reciprocal gamma products,
reflection and duplication, and the Mellin identification with the zeta
Dirichlet series. Its finite xi approximation has exact algebraic symmetry;
its identification with the analytic zeta function requires the substantive
Mellin proof. The archive explicitly says that its Lean interfaces are
uncompiled and incomplete. We reuse its written estimates at that stated
level, not as a kernel-checked theorem.

The Cauchy--residue draft gives finite triangulation cancellation,
quantitative Laurent expansions and rational contour winding calculations.
Its core singularity locations are supplied rational centres. A generic
argument principle for unknown zeta zeros does not follow merely by citing
that theorem: the local factorization, counting and approximation work
must also be supplied. Similarly, the Stirling draft proves the real
factorial formula with explicit remainders; it is not a proof of complex
log-gamma asymptotics. The complex estimate used in zero counting needs its
own proof. The FTA draft provides a written quantitative construction of
primitive recursive root names, but also explicitly lacks a completed
object-language Lean derivation.

\subsection{A concrete completion order}
The accompanying \texttt{ObjectTargets.lean} is an uncompiled interface.
It adds no analytic axiom and supplies no instance of its final obligation
structure. The implementation sequence is as follows.
\begin{enumerate}
\item Compile integer, rational, list, bounded-search and polynomial
      operations to the repository's \texttt{PR} syntax. Prove defining
      equations and the required arithmetic lemmas inside
      \texttt{Prf PRAAx}; semantic Lean calculations alone do not suffice.
\item Compile the fixed elementary and zeta/xi approximation programs.
      Formalize their rational error budgets, input sensitivities and
      derivative bounds with the displayed quantifier-free invariants.
\item Prove the finite functional-equation, Euler-product, and gamma
      identities. Keep numerical division guarded by an explicit positive
      lower bound. Prove the positive nonvanishing estimates.
\item Implement rational contour certification and winding. Prove
      finite homotopy invariance and the counting statements, including
      the multiplicity convention.
\item Formalize the quantitative Jensen, minimum-modulus, complex
      Stirling and argument-variation estimates. Compile the bounded
      regular-height generator and prove that its output passes the
      certificate checker.
\item Construct the internal derivations of the rational zero-counting
      and reciprocal-tail sentences. Derive unboundedness and the stated
      distribution corollaries from these internal finite estimates.
\item Apply \texttt{pra\_iff\_native} to each completed derivation.
      Kernel-check all resulting proof terms at the pinned toolchain,
      and audit dependencies for \texttt{sorryAx} and additional axioms.
\end{enumerate}
The inspected repository requests Lean and mathlib \texttt{v4.32.0}.
Those versions identify a proposed build environment; this package does not
claim that its new analytic development has been compiled there. A proof
about \texttt{Complex} in ambient Lean can be used as an independent
semantic comparison, but does not establish the restriction to iPRA.

\section{Proof status, archive contents, and reproducibility}
\label{sec:status}

\subsection{The three evidence levels}
The mathematical deliverable is the human argument in this report,
including its numerical bounds and finite-arithmetic interpretation.
The functional-equation foundation is reused as a written proof from
the included upstream archive. The uniform rational-polynomial and
winding arguments supply the finite infrastructure required by the
counting proof; no fixed-degree FTA schema is silently substituted for
the variable-degree construction.

The executable evidence consists of exact rational comparisons for the
worked simple zero and exact integer calculations of the regular-height
bounds. The latter program computes cutoff sizes and verifies their
elementary inequalities; it does not execute the enormous general
height search. A finite run is evidence about that run. The uniform
mathematical proof is the reason the algorithm has the stated behavior
on all inputs.

The Lean deliverable is a concrete interface and a module-by-module
completion plan. No new analytic Lean theorem has been kernel-checked
in this package. The interface has no \code{sorry} proof and adds no
analytic axiom, but no instance of its structure of final obligations
is constructed. A specification without a supplied proof is not a proof. The essential
completion target is an actual term of
\code{Prf PRAAx [] phi}; only then does
\code{pra_iff_native} transfer it to the concise native calculus.

\subsection{Files}
\begin{center}
\begin{tabular}{p{.43\textwidth}p{.49\textwidth}}
\toprule
Path & Purpose\\
\midrule
\code{zeta_zeros_ipra.tex}, \code{sections/}
 & Complete source of this report.\\
\code{zeta_zeros_ipra.pdf}
 & Compiled, visually inspected report.\\
\code{lean/ObjectTargets.lean}
 & Object-language sentence constructors, the verified repository bridge
   as an interface, and explicit final proof obligations.\\
\code{lean/README.md}
 & Proposed build location, dependency versions and completion gates.\\
\code{code/rational_zeta.py}
 & Exact rational zeta enclosures and the simple-zero certificate.\\
\code{code/good_height_bounds.py}
 & Exact integer size calculations for the theoretical regular-height
   construction.\\
\code{verification/}
 & Actual computation outputs and verification notes.\\
\code{sources/}
 & Pinned repository source snapshots and provenance manifest.\\
\path{dependencies/zeta_functional_equation_ipra.zip}
 & Unchanged upstream functional-equation proof, finite contracts and
   explicitly incomplete Lean outline.\\
\code{SHA256SUMS}
 & Hashes for the files in this archive.\\
\bottomrule
\end{tabular}
\end{center}

\subsection{Reproduction}
With Python 3.12 or a compatible later Python, the two scripts require
only the standard library. The main report uses ordinary LaTeX packages.
Run the supplied \code{build.sh} from the archive root to regenerate
the computation outputs and compile the PDF. The exact build used for
this delivery and the observed results are recorded in
\code{verification/STATUS.md}.

The repository snapshot is read-only in this work. Source files copied
from it were checked against their Git blob hashes, and their SHA-256
hashes are recorded. The included functional-equation ZIP has SHA-256
\begin{center}\small\ttfamily
78ddf1de0c4c5fc3e6a553efdeda96e617a646f7a02e62f84481d57c8691e4de.
\end{center}
No claim about newer repository revisions is needed for the argument.


\begin{thebibliography}{99}
\bibitem{MathWorld}
E. W. Weisstein, \emph{Riemann Zeta Function Zeros}, Wolfram MathWorld.
\url{https://mathworld.wolfram.com/RiemannZetaFunctionZeros.html}.
Consulted 15 November 2026; used to identify the requested subject.

\bibitem{DLMFreflection}
NIST Digital Library of Mathematical Functions, \S25.4,
\emph{Reflection Formulas}.
\url{https://dlmf.nist.gov/25.4}.

\bibitem{DLMFzeros}
NIST Digital Library of Mathematical Functions, \S25.10,
\emph{Zeros}.
\url{https://dlmf.nist.gov/25.10}.

\bibitem{DLMFgamma}
NIST Digital Library of Mathematical Functions, \S5.11,
\emph{Asymptotic Expansions}, and \S5.9, \emph{Integral Representations}.
\url{https://dlmf.nist.gov/5.11}; \url{https://dlmf.nist.gov/5.9}.

\bibitem{Backlund}
R. J. Backlund, \emph{\"Uber die Nullstellen der Riemannschen Zetafunktion},
Acta Mathematica \textbf{41}, 345--375 (1918 printing).
\url{https://archive.ymsc.tsinghua.edu.cn/pacm_download/117/5290-11511_2006_Article_BF02422950.pdf}.
The argument-variation and Jensen strategy is classical; the bounds and
bounded contour selection here are worked out explicitly.

\bibitem{FLINT}
FLINT documentation, \emph{Zeta functions, polylogarithms and L-functions}.
\url{https://flintlib.org/doc/algorithms/special-functions/zeta.html}.
Consulted for the Euler--Maclaurin representation; the remainder estimate
used in the executable is derived in this report.

\bibitem{Repo}
V. Reshetnikov, \emph{ProveIt}, snapshot
\code{b5a3ad037723},
\path{Logic/PrimitiveRecursiveArithmetic/ConciseIPRA}.
\url{https://github.com/VladimirReshetnikov/ProveIt/tree/b5a3ad037723c44d69a3d5b28a950c554327a989/Logic/PrimitiveRecursiveArithmetic/ConciseIPRA}.

\bibitem{RepoFE}
\emph{Riemann's functional equation in intuitionistic primitive recursive
arithmetic}, archived research draft in ProveIt,
\code{docs/drafts/zeta_functional_equation_ipra.zip}, at the snapshot
in~\cite{Repo}. Human proof and finite contracts; Lean development
explicitly incomplete.
\url{https://github.com/VladimirReshetnikov/ProveIt/blob/b5a3ad037723c44d69a3d5b28a950c554327a989/docs/drafts/zeta_functional_equation_ipra.zip}.

\bibitem{RepoAnalytic}
ProveIt archived research drafts,
\code{cauchy_residue_ipra.zip}, \code{stirling_ipra.zip}, and
\code{FTA_iPRA.zip}, at the snapshot in~\cite{Repo}.
Inspected for scope and proof status; not used as a blanket formal
analytic library. See the source manifest and
Section~\ref{sec:foundation-audit}.
\end{thebibliography}
\end{document}
